% generated by qpp/qpp2latex.py from qpp/examples/vocabulary.qpp; do not edit
% qpp: maps Loss Map Line Halves Dyads Values Mesh PosLine StepLine PosValues
% qpp: textrel _ zsite _
% qpp: textrel _ zcount _
% qpp: textrel _ zfloor _ zbox _
% qpp: textrel _ asymp _ on _
% qpp: textrel _ bdd _
% qpp: textrel _ below _ on _
% qpp: textrel _ zcap _
% qpp: textrel _ zcaps _ zby _
% qpp: textrel _ zfeasible _ zlevel _
% qpp: textrel _ zrows _
% qpp: textrel _ zweight _ zcorners _ zsides _
% qpp: textrel _ zpos _
% qpp: textrel _ zcapped _ zon _
% qpp: textrel _ zbox _ zin _
% qpp: textrel _ zroom _
% qpp: textrel _ zlvl _
% qpp: textrel _ zwide _
% qpp: textrel _ zfit _
% qpp: textrel _ zplane _
% qpp: relations coef coefw

The vocabulary rows over Mathlib, each used once: words, defines, relations, aliases, types, contexts and
typed words, with the proof words that read them.

\section{The vocabulary rows}\label{sec:vocabulary}

\begin{fact}[{the archimedean property, Mathlib \texttt{exists\_\allowbreak nat\_\allowbreak gt}}]\label{fact:arch}
The words of base, as base has them.
\begin{stmt}
\sit{core}{x\in \Q:\ \exists n\in \N:\ n > x}
\end{stmt}
\end{fact}

\begin{stmt}
\sit{word}{\text{mass}:=\Q_{>0}}
\sit{word}{\text{integer}:=\Z}
\sit{define}{\mathrm{pos}(x):=\max(x,0)}
\sit{alias}{\mathrm{Arch}:=\cref{fact:arch}}
\sit{word}{\text{weightset}:=\Q_{\ge0}}
\sit{define}{\mathrm{half}(x):=x/2}
\sit{define}{\mathrm{mean}(x,y):=(x+y)/2}
\sit{relation}{q\ \text{zsite}\ t:\iff q\in \Q_{>0},\ q \le t,\ t \le 2 q}
\sit{alias}{\mathrm{Zarch}:=\cref{fact:arch}}
\end{stmt}

\begin{lemma}[{a word for a set}]\label{lem:zword}
A word for a set.
\begin{stmt}
\sit{given}{x\in \mathrm{mass},\ y\in \mathrm{weightset},\ x \le y}
\sit{(1)}{x x \le x y}
\end{stmt}
\end{lemma}

\begin{proof}
\begin{plines}
\pline{(1)}{The goal follows from \(x\) times \((\mathrm{H}_{2})\).}
\end{plines}
\end{proof}

\begin{lemma}[{the fact of a step in a certificate}]\label{lem:zwithcert}
A certificate that names the fact of a step: hwK is the fact of the K-th step with of the proof line.
\begin{stmt}
\sit{given}{x,y\in \Q_{>0},\ x \le y}
\sit{(1)}{2 x \le x+y}
\end{stmt}
\end{lemma}

\begin{proof}
\begin{plines}
\pline{(1)}{\textup{(1)}~\(x+x \le x+y\) from \((\mathrm{H}_{1})\). The goal follows from \textup{(1)}.}
\end{plines}
\end{proof}

\begin{lemma}[{masses as multipliers in a certificate}]\label{lem:zmasscert}
A certificate with masses as multipliers: each multiplier, and each monomial of the rest, positive by the name of
its sign.
\begin{stmt}
\sit{given}{x,y,a,b\in \Q_{>0},\ k\in \N_{\ge2},\ w_1,\dots,w_k\in \Q_{>0},\ a \le b,\ x \le y}
\sit{(1)}{x a+x x \le x b+x y+2 x y}
\sit{(2)}{x < y+x y}
\sit{(3)}{x a < x b+k w_k y}
\sit{(4)}{2 a+x/2 \le 2 b+y/2}
\end{stmt}
\end{lemma}

\begin{proof}
\begin{plines}
\pline{(1)}{The goal follows combining \(x\) times \((\mathrm{H}_{1})\) and \(x\) times \((\mathrm{H}_{2})\).}
\pline{(2)}{The goal follows from \((\mathrm{H}_{2})\).}
\pline{(3)}{The goal follows from \(x\) times \((\mathrm{H}_{1})\).}
\pline{(4)}{The goal follows combining \(2\) times \((\mathrm{H}_{1})\) and \(1/2\) times \((\mathrm{H}_{2})\).}
\end{plines}
\end{proof}

\begin{lemma}[{equations in a certificate}]\label{lem:zeqcert}
Equations in a certificate: a relation between counts, read in QQ by its cast; an equation as the goal; a goal true
by ring, from a datum times 0; an equation read from its other side, after the bar of V[written | turned].
\begin{stmt}
\sit{given}{c,\mathrm{ca},\mathrm{cb}\in \N,\ x,y,z\in \Q_{>0},\ c = \mathrm{ca}+\mathrm{cb},\ x \le c,\ y = x+z}
\sit{(1)}{x \le \mathrm{ca}+\mathrm{cb}}
\sit{(2)}{y+x = 2 x+z}
\sit{(3)}{0 x y \le 0}
\sit{(4)}{x+z \le y}
\end{stmt}
\end{lemma}

\begin{proof}
\begin{plines}
\pline{(1)}{The goal follows combining \((\mathrm{H}_{1})\) and \((\mathrm{H}_{2})\).}
\pline{(2)}{The goal follows from \((\mathrm{H}_{3})\).}
\pline{(3)}{The goal follows from \(0\) times \((\mathrm{H}_{2})\).}
\pline{(4)}{The goal follows from \((\mathrm{H}_{3})\) read from right to left.}
\end{plines}
\end{proof}

\begin{lemma}[{case facts, instances and no fact in a certificate}]\label{lem:zcasecert}
The data of V also name the clause of a case split and its negation (hcK, hnK), a universal fact at an index in
scope (h_1(i)), or no fact at all (V[0]: the rest, a sum of positive monomials, closes the goal).
\begin{stmt}
\sit{index}{n\in \N,\ a\in \mathrm{Point}(n),\ x\in \Q_{>0},\ (\forall i:\ a_i \le x):\ \forall i:\ a_i \le 2 x}
\sit{rest}{x\in \Q_{>0}:\ x x \le x x+x}
\sit{order}{x,y\in \Q_{\ge0},\ (y < x\Rightarrow x \le y+1):\ x \le y+1}
\end{stmt}
\end{lemma}

\begin{proof}
\begin{plines}
\pline{index}{The goal follows from \((\mathrm{H}_{1})\) at \(i\).}
\pline{rest}{The goal follows by computation.}
\pline{order}{Case \((\mathrm{c}_{1})\) \(x \le y\): the goal follows from \((\mathrm{c}_{1})\). Case \((\bar{\mathrm{c}}_{1})\) \(x > y\): the goal follows by logic and from the facts at hand.}
\end{plines}
\end{proof}

\begin{lemma}[{a define}]\label{lem:zdef}
A define, as a tie and under cases.
\begin{stmt}
\sit{given}{x,y\in \Q_{>0},\ x \le y}
\sit{(1)}{\mathrm{half}(x) \le \mathrm{half}(y)}
\sit{(2)}{\mathrm{pos}(x-y) = 0}
\sit{(3)}{\mathrm{pos}(y-x)+x = y}
\sit{(4)}{\max(x,y,2 y) = 2 y}
\sit{(5)}{\mathrm{mean}(x,y) \le y}
\sit{(6)}{\min(x,y,2 y) = x}
\sit{(7)}{\max(y-x,0)+x = y}
\sit{(8)}{\max(x-y,0) = 0}
\end{stmt}
\end{lemma}

\begin{proof}
\begin{plines}
\pline{(1)}{The goal follows from the facts at hand.}
\pline{(2)}{The goal follows by cases on the conditionals of the facts, then by computation in each case.}
\pline{(3)}{The goal follows by cases on the conditionals of the facts, then by computation in each case.}
\pline{(4)}{The goal follows by cases on the conditionals of the facts, then by computation in each case.}
\pline{(5)}{The goal follows from the facts at hand.}
\pline{(6)}{The goal follows by cases on the conditionals of the facts, then by computation in each case.}
\pline{(7)}{The goal follows by cases on the conditionals of the facts, then by computation in each case.}
\pline{(8)}{The goal follows by cases on the conditionals of the facts, then by computation in each case.}
\end{plines}
\end{proof}

\begin{lemma}[{a relation}]\label{lem:zrel}
A relation as one hypothesis, its rows read by projection.
\begin{stmt}
\sit{given}{q,t\in \Q,\ q\ \text{zsite}\ t}
\sit{(1)}{t \le 2 q}
\sit{(2)}{q q \le q t}
\sit{(3)}{3 q \le 4 t}
\end{stmt}
\end{lemma}

\begin{proof}
\begin{plines}
\pline{(1)}{The goal follows from the facts at hand.}
\pline{(2)}{The goal follows from the facts at hand.}
\pline{(3)}{The goal follows from the facts at hand.}
\end{plines}
\end{proof}

\begin{lemma}[{a division under a relation, the sign in the binder}]\label{lem:zreldivpos}
The same with the sign in the binder.
\begin{stmt}
\sit{given}{q,t\in \Q_{>0},\ q\ \text{zsite}\ t}
\sit{(1)}{t/q \le 2}
\end{stmt}
\end{lemma}

\begin{proof}
\begin{plines}
\pline{(1)}{The goal follows dividing by \(q\).}
\end{plines}
\end{proof}

\begin{lemma}[{mono}]\label{lem:zmono}
A monotone comparison.
\begin{stmt}
\sit{given}{a,b,c\in \Q_{>0},\ a \le b}
\sit{(1)}{a c+a^2 \le b c+b^2}
\end{stmt}
\end{lemma}

\begin{proof}
\begin{plines}
\pline{(1)}{The goal follows by monotonicity.}
\end{plines}
\end{proof}

\begin{lemma}[{mono on powers of one base}]\label{lem:zmonopow}
A power of one base on both sides: a base at least 1 keeps the order of the exponents, a base at most 1 reverses
it.
\begin{stmt}
\sit{given}{y,z\in \Q_{\ge0},\ \eta\in \N,\ 1 \le y,\ z \le 1}
\sit{(1)}{y^2 \le y^{3+\eta}}
\sit{(2)}{z^{3+\eta} \le z^2}
\end{stmt}
\end{lemma}

\begin{proof}
\begin{plines}
\pline{(1)}{The goal follows by monotonicity.}
\pline{(2)}{The goal follows by monotonicity.}
\end{plines}
\end{proof}

\begin{lemma}[{an equation from two bounds}]\label{lem:zantisymm}
An equation of rationals from its two bounds (le_antisymm).
\begin{stmt}
\sit{given}{a,b\in \Q_{\ge0},\ a \le b,\ b \le a}
\sit{(1)}{a = b}
\end{stmt}
\end{lemma}

\begin{proof}
\begin{plines}
\pline{(1)}{Each part of the goal follows from the facts at hand.}
\end{plines}
\end{proof}

\begin{lemma}[{a bound on the binder}]\label{lem:natbound}
A count with its bound on the binder (k in NN>=3): linarith reads the bound cast to the rationals, as it reads a
hypothesis 3 <= k of the given.
\begin{stmt}
\sit{given}{k\in \N_{\ge3},\ x\in \Q_{>0}}
\sit{(1)}{3 x \le k x}
\end{stmt}
\end{lemma}

\begin{proof}
\begin{plines}
\pline{(1)}{The goal follows by monotonicity.}
\end{plines}
\end{proof}

\begin{lemma}[{core items}]\label{lem:coreitems}
Core items: each declares its names and is read without the given (none of them has x <= 1); `core` and
`core NAME` are read by the checks only, `core 1` is printed.
\begin{stmt}
\sit{given}{x\in \Q_{>0},\ x \le 1}
\sit{(1)}{x x \le x}
\sit{core 1}{x\in \Q_{\ge0}:\ x \le x+x}
\end{stmt}
\end{lemma}

\begin{proof}
\begin{plines}
\pline{(1)}{The goal follows from the facts at hand.}
\pline{core 1}{The goal follows from the facts at hand.}
\end{plines}
\end{proof}

\begin{fact}[{a sum bounded term by term, Mathlib \texttt{Finset.\allowbreak sum\_\allowbreak le\_\allowbreak sum}}]\label{fact:termwise}
Sums by cited rules of Mathlib: a family given by a formula in its index (`b_(i) := c`), a side goal of the cited
rule proved by its word (`HYP by W`), and sum[W] with the hypotheses over the range at the index.
\end{fact}

\begin{fact}[{a sum of equal terms, Mathlib \texttt{Finset.\allowbreak sum\_\allowbreak const}}]\label{fact:constsum}
\begin{stmt}
\sit{core}{N\in \N,\ c\in \Q_{\ge0}:\ \textstyle\sum_{i=1}^N c = N c}
\end{stmt}
\end{fact}

\begin{fact}[{the pigeonhole principle for sums, Mathlib \texttt{Finset.\allowbreak exists\_\allowbreak lt\_\allowbreak of\_\allowbreak sum\_\allowbreak lt}}]\label{fact:pigeon}
\begin{stmt}
\sit{core}{N\in \N,\ f_1,\dots,f_N,g_1,\dots,g_N\in \Q_{\ge0},\ \textstyle\sum_{i=1}^N f_i < \textstyle\sum_{i=1}^N g_i:\ \exists i\in \N_{\ge1}:\ i \le N,\ f_i < g_i}
\end{stmt}
\end{fact}

\begin{fact}[{a gap between two counts, Mathlib \texttt{lt\_\allowbreak iff\_\allowbreak exists\_\allowbreak add.\allowbreak mp}}]\label{fact:zgap}
A cited rule that takes the premise of the statement as its argument, the guards of its names left out: the rule
applied to the premise by its name.
\end{fact}

\begin{lemma}[{terms at most c: the sum at most N c, term by term}]\label{lem:sumupper}
\end{lemma}

\begin{lemma}[{the same by the cited rule at the constant family}]\label{lem:sumupperref}
\end{lemma}

\begin{lemma}[{the same with the universal side goal written and proved by V}]\label{lem:sumupperside}
A side goal of the cited rule quantified over the index range, proved by V with its certificate: at each index j of
the range it is the hypothesis h_3 of the row at j, whose premises are those of the side goal.
\end{lemma}

\begin{lemma}[{a sum above N G has a term above G}]\label{lem:sumscale}
\end{lemma}

\begin{lemma}[{a sum above N G_i has a term above G_i}]\label{lem:sumscalefam}
The same with the data a member of a family, the side goal in the names of the row (a name free in the data is
never captured by a name the rule binds), and printed.
\begin{stmt}
\sit{given}{N\in \N,\ i\in \N_{\ge1},\ g_1,\dots,g_N,G_1,\dots,G_N\in \Q_{\ge0},\ \textstyle\sum_{j=1}^N g_j > N G_i}
\sit{(1)}{\exists j\in \N_{\ge1}:\ j \le N,\ G_i < g_j}
\end{stmt}
\end{lemma}

\begin{proof}
\begin{plines}
\pline{(1)}{The goal follows by \cref{fact:pigeon} (\(\textstyle\sum_{j=1}^N G_i < \textstyle\sum_{j=1}^N g_j\) by \cref{fact:constsum}).}
\end{plines}
\end{proof}

\begin{lemma}[{sums of binomials against a geometric weight}]\label{lem:zpascal}
The word pascal: the sums of binomials against the weight 1/2^j are at most 2^(k + 1), by induction on k for every
J (Pascal's rule shifted, S_(k+1)(J+1) = S_k(J+1) + S_(k+1)(J)/2, and the absorption of the last term, the terms
being at least 0); and the sums of the powers (j + 1)^k against the same weight, from (j + 1)^k <= k! binom(j + k, k).
\begin{stmt}
\sit{given}{k,J\in \N}
\sit{(1)}{\textstyle\sum_{j=0}^J \tbinom{j+k}{k}/2^j \le 2^{k+1}}
\sit{(2)}{\textstyle\sum_{j=0}^J (j+1)^k/2^j \le (\textstyle\prod_{i=1}^k i)\cdot 2^{k+1}}
\end{stmt}
\end{lemma}

\begin{proof}
\begin{plines}
\pline{(1)}{The goal follows by the bound of the binomial sums.}
\pline{(2)}{The goal follows by the bound of the binomial sums.}
\end{plines}
\end{proof}

\begin{lemma}[{an alias}]\label{lem:zarch}
An alias of this module and one of base.
\begin{stmt}
\sit{given}{x\in \Q}
\sit{(1)}{\exists n\in \N:\ n > x}
\sit{(2)}{\exists n\in \N:\ n > x}
\end{stmt}
\end{lemma}

\begin{proof}
\begin{plines}
\pline{(1)}{The goal follows by \cref{fact:arch}.}
\pline{(2)}{The goal follows by \cref{fact:arch}.}
\end{plines}
\end{proof}

\begin{lemma}[{cases with data}]\label{lem:zcases}
Cases with data.
\begin{stmt}
\sit{given}{x,y\in \Q_{\ge0}}
\sit{(1)}{x y \le \max(x,y) \max(x,y)}
\end{stmt}
\end{lemma}

\begin{proof}
\begin{plines}
\pline{(1)}{The goal follows by cases on the conditionals of the facts, then from the product of \((\mathrm{H}_{1})\) and \((\mathrm{H}_{2})\) in each case.}
\end{plines}
\end{proof}

\begin{lemma}[{a strict division}]\label{lem:zstrict}
A strict division.
\begin{stmt}
\sit{given}{a,b,c\in \Q_{>0},\ c a > c b}
\sit{(1)}{a > b}
\end{stmt}
\end{lemma}

\begin{proof}
\begin{plines}
\pline{(1)}{The goal follows dividing by \(c\).}
\end{plines}
\end{proof}

\begin{lemma}[{a double minus}]\label{lem:zneg}
A double minus.
\begin{stmt}
\sit{given}{t\in \Q}
\sit{(1)}{-(-t) = t}
\end{stmt}
\end{lemma}

\begin{proof}
\begin{plines}
\pline{(1)}{The goal follows by substitution.}
\end{plines}
\end{proof}

\begin{stmt}
\sit{define}{\mathrm{ordz}(m):=m-1}
\sit{relation}{m\ \text{zcount}\ t:\iff m\in \N_{\ge1},\ t\in \Q_{\ge0},\ m \le t+1}
\end{stmt}

\begin{lemma}[{a relation with a count}]\label{lem:znat}
A relation with a count, and a define applied to a count.
\begin{stmt}
\sit{given}{m\in \N,\ t\in \Q,\ m\ \text{zcount}\ t}
\sit{(1)}{\mathrm{ordz}(m) \le t}
\sit{(2)}{m \ge 1}
\sit{(3)}{\mathrm{ordz}(m) \ge 0}
\end{stmt}
\end{lemma}

\begin{proof}
\begin{plines}
\pline{(1)}{The goal follows from the facts at hand.}
\pline{(2)}{The goal follows from the facts at hand.}
\pline{(3)}{The goal follows from the facts at hand.}
\end{plines}
\end{proof}

\begin{lemma}[{a relation between two terms}]\label{lem:znat2}
The relation applied to the count m + 1: the cast of the sum in its row is pushed.
\begin{stmt}
\sit{given}{a,b\in \Q_{>0},\ m\in \N_{\ge1},\ m+1\ \text{zcount}\ a+b}
\sit{(1)}{m \le a+b}
\end{stmt}
\end{lemma}

\begin{proof}
\begin{plines}
\pline{(1)}{The goal follows from the facts at hand.}
\end{plines}
\end{proof}

\begin{lemma}[{a division under a relation}]\label{lem:zreldiv}
\begin{stmt}
\sit{given}{q,t\in \Q,\ q\ \text{zsite}\ t}
\sit{(1)}{t/q \le 2}
\end{stmt}
\end{lemma}

\begin{proof}
\begin{plines}
\pline{(1)}{The goal follows dividing by \(q\).}
\end{plines}
\end{proof}

\begin{lemma}[{a certificate that names a row of the relation}]\label{lem:zdata}
A certificate that names a row of a relation (h_1_2: the second row of h_1).
\begin{stmt}
\sit{given}{q,t\in \Q,\ q\ \text{zsite}\ t}
\sit{(1)}{q q \le q t}
\end{stmt}
\end{lemma}

\begin{proof}
\begin{plines}
\pline{(1)}{The goal follows from \(q\) times \((\mathrm{H}_{1.2})\).}
\end{plines}
\end{proof}

\subsection{Types, contexts and typed words}\label{sec:types}

Types and words are built from the primitives of Q++ only: tuples X^n, functions (X -> Y), sets set(X), pair;
join(x, y) concatenates two tuples and B times C is the product of two sets of tuples. No row names Lean or
Mathlib.
two halves of one size: the equation is proved by the lemmas of Fin.append only (simp only)
A word with an argument in NN (a depth N, which its result type may name), whose values are tuples: defined at a
coordinate, halve(n, N)(u)_(j); its body takes powers, sums over a range and brackets of natural numbers.
A relation on a tuple of tuples, read entry by entry (one subscript in print, f_{\iota,i}), and a relation with another
one as a clause.
A bundle: a program as one object, its parts named Pr.P, ...; a relation on it, a clause wherever it is used,
inside a bracket too.
A bundle with a part in NN, read part by part: the tuple by its coordinates, the count as a natural number.

\begin{stmt}
\sit{type}{\mathrm{Pair}:=(\Q_{\ge0}\times\N_{\ge1})\cup\{\top\}}
\sit{type}{\mathrm{Point}(n):=(\Q_{\ge0})^{n}}
\sit{type}{\mathrm{Loss}(n):=\mathrm{Point}(n)\to \Q_{\ge0}}
\sit{type}{\mathrm{Weight}(n):=(\Q_{>0})^{n}}
\sit{type}{\mathrm{Region}(n):=\mathcal{P}(\mathrm{Point}(n))}
\sit{type}{\mathrm{Map}(n):=\mathrm{Point}(n)\to \mathrm{Point}(n)}
\sit{type}{\mathrm{Grid}(n):=(\N)^{n}}
\sit{type}{\mathrm{Rows}(k,n):=((\N)^{n})^{k}}
\sit{type}{\mathrm{Forms}(k,n):=((\Q_{\ge0})^{n})^{k}}
\sit{type}{\mathrm{Sides}(k,n):=((\Q_{>0})^{n})^{k}}
\sit{type}{\mathrm{Line}:=\N\to \Q_{\ge0}}
\sit{word}{\mathrm{sep}(K\in \mathrm{Loss}(n),\ K'\in \mathrm{Loss}(n'))\in \mathrm{Loss}(n+n')}
\sit{define}{\mathrm{sep}(K,K')(x,y):=K(x)+K'(y)}
\sit{word}{\mathrm{shift}(K\in \mathrm{Loss}(n))\in \mathrm{Loss}(n)}
\sit{define}{\mathrm{shift}(K)(w):=2 K(w)+1}
\sit{word}{\mathrm{twice}(K\in \mathrm{Loss}(n),\ K'\in \mathrm{Loss}(n))\in \mathrm{Loss}(n+n)}
\sit{define}{\mathrm{twice}(K,K')(x,y):=K(x) K'(y)}
\sit{type}{\mathrm{Halves}(n):=\mathrm{Grid}(n)\to \mathrm{Point}(n)}
\sit{type}{\mathrm{Dyads}(n):=\mathrm{Grid}(n)\to \mathrm{Grid}(n)}
\sit{type}{\mathrm{Values}(n):=\mathrm{Grid}(n)\to \Q_{\ge0}}
\sit{type}{\mathrm{Depth}(n):=\{F\in \mathrm{Halves}(n),\ N\in \N\}}
\sit{word}{\mathrm{halve}(n,N\in \N)\in \mathrm{Halves}(n)}
\sit{define}{\mathrm{halve}(n,N)(u)_{j}:=u_j/2^N}
\sit{word}{\mathrm{dyad}(B\in \mathrm{Grid}(n),\ N\in \N)\in \mathrm{Dyads}(n)}
\sit{define}{\mathrm{dyad}(B,N)(u)_{j}:=\textstyle\sum_{k=1}^{B_j} \mathbf 1[u_j\cdot 2^k \le 2^N]}
\sit{word}{\mathrm{value}(K\in \mathrm{Loss}(n),\ N\in \N)\in \mathrm{Values}(n)}
\sit{define}{\mathrm{value}(K,N)(u):=K(\mathrm{halve}(n,N)(u))}
\sit{relation}{G\ \text{zfloor}\ Q\ \text{zbox}\ b:\iff n\in \N,\ G\in \mathrm{Values}(n),\ Q\in \Q_{>0},\ b\in \mathrm{Grid}(n),\ \forall u\in \mathrm{Grid}(n):\ (\forall i:\ u_i \le b_i)\Rightarrow Q \le G(u)}
\sit{given relation}{\mathrm{coef}(K\in \mathrm{Loss}(n),\ B\in \mathrm{Region}(n),\ p\in \mathrm{Pair})}
\sit{given relation}{\mathrm{coefw}(K\in \mathrm{Loss}(n),\ \omega\in \mathrm{Weight}(n),\ B\in \mathrm{Region}(n),\ p\in \mathrm{Pair})}
\sit{relation}{K\ \text{asymp}\ K'\ \text{on}\ B:\iff n\in \N,\ K,K'\in \mathrm{Loss}(n),\ B\in \mathrm{Region}(n),\ \exists c,C\in \Q_{>0}:\ \forall w\in B:\ c K'(w) \le K(w) \le C K'(w)}
\sit{relation}{K\ \text{bdd}\ B:\iff n\in \N,\ K\in \mathrm{Loss}(n),\ B\in \mathrm{Region}(n),\ \exists c\in \Q_{>0}:\ \forall w\in B:\ c \le K(w)}
\sit{relation}{K\ \text{below}\ K'\ \text{on}\ B:\iff n\in \N,\ K,K'\in \mathrm{Loss}(n),\ B\in \mathrm{Region}(n),\ \exists C\in \Q_{>0}:\ \forall w\in B:\ K(w) \le C K'(w)}
\sit{relation}{f\ \text{zcap}\ t:\iff k,n\in \N,\ f\in \mathrm{Rows}(k,n),\ t\in \N,\ \forall \iota:\ \forall i:\ f_{\iota,i} \le t}
\sit{relation}{f\ \text{zcaps}\ g\ \text{zby}\ t:\iff k,n,t\in \N,\ f,g\in \mathrm{Rows}(k,n),\ f\ \text{zcap}\ t,\ \forall j:\ \forall i:\ g_{ji} \le f_{ji}}
\sit{type}{\mathrm{Prog}(k,n):=\{P,N\in \mathrm{Rows}(k,n),\ \mathrm{op},\mathrm{om},\mathrm{eps}\in \mathrm{Grid}(k)\}}
\sit{relation}{s\ \text{zfeasible}\ \mathrm{Pr}\ \text{zlevel}\ L:\iff k,n,L\in \N,\ \mathrm{Pr}\in \mathrm{Prog}(k,n),\ s\in \mathrm{Grid}(n),\ \forall \iota:\ \textstyle\sum_i \mathrm{Pr}.P_{\iota,i}\cdot s_i+\mathrm{Pr}.\mathrm{op}_{\iota} \ge \textstyle\sum_i \mathrm{Pr}.N_{\iota,i}\cdot s_i+\mathrm{Pr}.\mathrm{om}_{\iota}+\mathrm{Pr}.\mathrm{eps}_{\iota}\cdot L}
\sit{type}{\mathrm{ZCert}(k):=\{y\in \mathrm{Grid}(k),\ d\in \N\}}
\end{stmt}

\begin{fact}[{Mono: a larger region has a smaller pair, \cite{doi_10_1017_CBO9780511800474}}]\label{fact:mono}
\end{fact}

\begin{fact}[{Min: on a union the pair is the least, \cite{doi_10_1017_CBO9780511800474}}]\label{fact:min}
\end{fact}

\begin{lemma}[{the pair is unique}]\label{lem:unique}
\end{lemma}

\begin{lemma}[{uses lem:unique}]\label{lem:uniq2}
\end{lemma}

\begin{lemma}[{the pair is unique, as an equation}]\label{lem:uniqeq}
An equation from two words, the first for <= and the second for >= (le_antisymm): by two lemmas, and by the
cited rule twice.
\end{lemma}

\begin{lemma}[{the pair is unique, from the cited rule twice}]\label{lem:uniqcit}
\end{lemma}

\begin{lemma}[{the product is monotone}]\label{lem:sum}
\end{lemma}

\begin{fact}[{Add: independent parts multiply, \cite{doi_10_1017_CBO9780511800474}}]\label{fact:add}
\end{fact}

\begin{fact}[{Min: on a union of pieces the pair is the least, \cite{doi_10_1017_CBO9780511800474}}]\label{fact:pieces}
\end{fact}

\begin{fact}[{a doi with the mark that closes a Lean comment, \cite{doi_10_1017_CBO9780511800474___exit}}]\label{fact:doimark}
A doi may hold the mark that closes a Lean comment: it is printed escaped in the doc comment of the cited fact,
so the build of this file checks that it closes nothing.
\end{fact}

\begin{lemma}[{a cited fact as a hypothesis}]\label{lem:zdoicit}
A cited fact of a doi with no context, a hypothesis of the row that cites it: over the kernel the row takes the
fact over the kernel, read in Mathlib by its reading.
\end{lemma}

\begin{fact}[{the empty region has the pair top, \cite{doi_10_1017_CBO9780511800474}}]\label{fact:empty}
\end{fact}

\begin{lemma}[{Add with the parts given}]\label{lem:addsym}
\end{lemma}

\begin{lemma}[{Min applied to the pieces}]\label{lem:pieces}
\end{lemma}

\begin{lemma}[{top is the unit of the tropical sum}]\label{lem:topunit}
\end{lemma}

\begin{lemma}[{top absorbs the product}]\label{lem:topabs}
\end{lemma}

\begin{lemma}[{the order is transitive}]\label{lem:trans}
\end{lemma}

\begin{lemma}[{the tropical sum is below each part}]\label{lem:minle}
\end{lemma}

\begin{lemma}[{the empty region is a unit of the union}]\label{lem:emptyunion}
\end{lemma}

\begin{lemma}[{the algebra of sets}]\label{lem:setalg}
\end{lemma}

\begin{lemma}[{the tropical sum is below each member}]\label{lem:infle}
\end{lemma}

\begin{lemma}[{a lower bound of every member is below the tropical sum}]\label{lem:leinf}
\end{lemma}

\begin{lemma}[{the least member is the tropical sum}]\label{lem:least}
\end{lemma}

\begin{fact}[{Chart: comparable losses have the same pair, \cite{doi_10_1017_CBO9780511800474}}]\label{fact:chart}
\end{fact}

\begin{fact}[{a loss bounded below has the pair top, \cite{doi_10_1017_CBO9780511800474}}]\label{fact:bdd}
\end{fact}

\begin{fact}[{Below: a smaller loss has a smaller pair, \cite[Proposition 3.4]{doi_10_18409_jas_v8i1_47}}]\label{fact:below}
\end{fact}

\begin{lemma}[{a loss is comparable to itself}]\label{lem:self}
\end{lemma}

\begin{lemma}[{Chart applied}]\label{lem:chart2}
\end{lemma}

\begin{lemma}[{a loss is nonnegative}]\label{lem:nonneg}
\end{lemma}

\begin{lemma}[{a loss after a map}]\label{lem:map}
\end{lemma}

\begin{lemma}[{a typed word at a point}]\label{lem:wpoint}
A typed word applied at a point: its value there is read through its equation NAME_def, a tie
for V; a word of two points at join(x, y), where its equation holds.
\end{lemma}

\begin{lemma}[{a word of two points at a joined point}]\label{lem:wjoin}
\end{lemma}

\begin{lemma}[{hypotheses over the points of a region, at a point}]\label{lem:winst}
A hypothesis quantified over the points of a region, used at the points the goal introduces: from
forall v in B : P and the point w of B, the instance P at w, a fact for V.
\end{lemma}

\begin{lemma}[{a hypothesis over points with a typed word}]\label{lem:winst2}
\end{lemma}

\begin{lemma}[{a word at a point in a hypothesis}]\label{lem:whyp}
\end{lemma}

\begin{lemma}[{a word with a depth, its values tuples}]\label{lem:wdepth}
A word with a depth, its values tuples: a statement names no word (it is not read over the kernel), and a proof
takes the word as a witness. V reads a coordinate of the word at a point through its equation, and through a
hypothesis that names its value there (y = halve(n, N)(u)); a relation takes a word as its argument; take gives a
part of a bundle by a word.
\end{lemma}

\subsection{Literal pairs}\label{sec:literals}

A tuple (l, m) is a literal pair where a type row `type NAME := pair` is in scope (a primitive of the list):
NAME.mk l (m - 1), the first coordinate built in NNRat from numerals, counts,
masses and names in QQ>=0 with + * / ^ min max, the second a multiplicity, at least 1 whatever its names (a
numeral at least 1, a name in NN>=k with k >= 1, + * ^ of those, a product over a range of those). A comparison of
sums of literals is read in Q with no branch: the word of the proof line compares the first coordinates, weakly
where the orders compare outright (omega), else strictly.

\begin{lemma}[{a product of literal pairs}]\label{lem:lsum}
\end{lemma}

\begin{lemma}[{a strict comparison of the first coordinates}]\label{lem:lstrict}
\end{lemma}

\begin{lemma}[{equal first coordinates: the orders compare}]\label{lem:lord}
\end{lemma}

\begin{lemma}[{the pair is monotone in its first coordinate}]\label{lem:lmono}
\end{lemma}

\begin{lemma}[{a min of first coordinates}]\label{lem:lmin}
\end{lemma}

\begin{lemma}[{orders that do not compare: the first coordinates compare strictly}]\label{lem:lorders}
\end{lemma}

\begin{lemma}[{a member of a family of masses}]\label{lem:lfam}
\end{lemma}

\begin{lemma}[{the product of three pairs}]\label{lem:lthree}
\end{lemma}

\begin{lemma}[{a literal pair is below top}]\label{lem:ltop}
\end{lemma}

\begin{fact}[{a cited rule with a literal pair, \cite{doi_10_1017_CBO9780511800474}}]\label{fact:half}
\end{fact}

\begin{lemma}[{a literal pair in a context}]\label{lem:lctx}
\end{lemma}

\begin{lemma}[{a name in QQ>=0}]\label{lem:lnn}
\end{lemma}

\begin{lemma}[{multiplicities: a power, a product over a range, a member of a family of counts}]\label{lem:lmult}
\end{lemma}

\begin{lemma}[{literal pairs over the given of the item}]\label{lem:lgiven}
\begin{stmt}
\sit{given}{x\in \Q_{>0},\ m\in \N_{\ge1}}
\sit{check}{(x,m+1) \le (x,m)}
\end{stmt}
\end{lemma}

\begin{proof}
\begin{plines}
\pline{check}{The goal follows from the facts at hand.}
\end{plines}
\end{proof}

\begin{lemma}[{a certificate on the first coordinates}]\label{lem:lcert}
\end{lemma}

A min of literal pairs (the tropical sum) is reduced with no search: under = its argument equal to the other side
whatever the names (min_eq_left or min_eq_right, the comparison of its two arguments as above); a min on the right
of <= or < by le_min or lt_min; a min on the left by the argument equal to the other side, or below it with no name
or by a hypothesis on the first coordinates.

\begin{lemma}[{a min of numeral pairs}]\label{lem:lminnum}
\end{lemma}

\begin{lemma}[{a min equal to one of its arguments whatever the names}]\label{lem:lminid}
\end{lemma}

\begin{lemma}[{a min equal to its argument by a hypothesis}]\label{lem:lminhyp}
\end{lemma}

\begin{lemma}[{below a min, and a min below}]\label{lem:lminle}
\end{lemma}

\begin{lemma}[{a min below by a hypothesis on the first coordinates}]\label{lem:lminlt}
\end{lemma}

A context whose pair is equal to the pair of a hypothesis whatever the names: `convert h using 1`, then the
equation of the two pairs as above.

\begin{lemma}[{a context pair equal by arithmetic}]\label{lem:lctxeq}
\end{lemma}

\begin{lemma}[{a context pair equal to a product of pairs}]\label{lem:lctxsum}
\end{lemma}

\begin{lemma}[{a context pair equal to a min}]\label{lem:lctxmin}
\end{lemma}

\subsection{Programs}\label{sec:programs}

The atom prog(DECLS, ROWS : COST) = P (a primitive of the list): the scales are
declared in QQ>=0, each row is a relation >=, <= or = of two affine forms in the scales, the cost is a linear form
in the scales, and P is a pair. prog(...) = (l, m) says that l is the least value of the cost over the
feasible points and that the points of cost l span an affine space of dimension m - 1; prog(...) = top says that
no point is feasible. The word prog[PRICES; POINTS] proves it: one price per row, then each point as assignments
(a scale left out is 0). The prices bound the value from below (weak duality) and the dimension from above (the
rows of positive price and the scales of positive reduced cost are tight on the face); the points attain the value
and bound the dimension from below. A sign of a price or of a reduced cost is read from the given; a fixture
computes the pair by vertex enumeration, without the word.

\begin{lemma}[{a program at a tie}]\label{lem:progtie}
One unit at the tie N = 2 alpha: the optimal face is the segment between the two vertices.
\begin{stmt}
\sit{given}{\alpha,N\in \Q_{>0},\ N = 2 \alpha}
\sit{(1)}{\operatorname{prog}(v,w\in \Q_{\ge0},\ 2 v+4 w \ge 1:\ \alpha v+N w) = (\alpha/2,2)}
\end{stmt}
\end{lemma}

\begin{proof}
Both reduced costs are zero at the tie (the second after N = 2 alpha is substituted): the priced row alone bounds
the face from above, and the two vertices bound it from below.
\begin{plines}
\pline{(1)}{The goal follows by the prices \(\alpha/2\) of the rows at the points \(v=1/2, w=0\); \(v=0, w=1/4\).}
\end{plines}
\end{proof}

\begin{lemma}[{a program above the tie}]\label{lem:progabove}
Off the tie, one item per case (the sign of a reduced cost is read from the given): the reduced cost N - 2 alpha of w
is positive by the hypothesis 2 alpha < N, so the face is the vertex v = 1/2.
\begin{stmt}
\sit{given}{\alpha,N\in \Q_{>0},\ 2 \alpha < N}
\sit{(1)}{\operatorname{prog}(v,w\in \Q_{\ge0},\ 2 v+4 w \ge 1:\ \alpha v+N w) = (\alpha/2,1)}
\end{stmt}
\end{lemma}

\begin{proof}
\begin{plines}
\pline{(1)}{The goal follows by the prices \(\alpha/2\) of the rows at the points \(v=1/2\).}
\end{plines}
\end{proof}

\begin{lemma}[{a program with a face of dimension 2}]\label{lem:progface}
An equation row and a face of dimension 2: the cost is constant on the triangle.
\begin{stmt}
\sit{given}{\alpha\in \Q_{>0}}
\sit{(1)}{\operatorname{prog}(a,b,c\in \Q_{\ge0},\ a+b+c = 1:\ \alpha a+\alpha b+\alpha c) = (\alpha,3)}
\end{stmt}
\end{lemma}

\begin{proof}
\begin{plines}
\pline{(1)}{The goal follows by the prices \(\alpha\) of the rows at the points \(a=1\); \(b=1\); \(c=1\).}
\end{plines}
\end{proof}

\begin{lemma}[{a program with a parametric price}]\label{lem:progden}
A parameter in a denominator: the price 1/(3 + eta) leaves the reduced cost (2 + eta)/(3 + eta) on w.
\begin{stmt}
\sit{given}{\eta\in \Q_{\ge0}}
\sit{(1)}{\operatorname{prog}(v,w\in \Q_{\ge0},\ (3+\eta) v+w \ge 1:\ v+w) = (1/(3+\eta),1)}
\end{stmt}
\end{lemma}

\begin{proof}
\begin{plines}
\pline{(1)}{The goal follows by the prices \(1/(3+\eta)\) of the rows at the points \(v=1/(3+\eta)\).}
\end{plines}
\end{proof}

\begin{lemma}[{a program whose face is a ray}]\label{lem:progray}
A scale absent from the cost: the face is a ray, and a point on it bounds the dimension from below.
\begin{stmt}
\sit{given}{\alpha\in \Q_{>0}}
\sit{(1)}{\operatorname{prog}(v,w\in \Q_{\ge0},\ v \ge 1:\ \alpha v) = (\alpha,2)}
\end{stmt}
\end{lemma}

\begin{proof}
\begin{plines}
\pline{(1)}{The goal follows by the prices \(\alpha\) of the rows at the points \(v=1\); \(v=1, w=1\).}
\end{plines}
\end{proof}

\begin{lemma}[{an empty program}]\label{lem:progempty}
No feasible point: the price 1 on the row gives the line 0 >= v + 1.
\begin{stmt}
\sit{given}{\alpha\in \Q_{>0}}
\sit{(1)}{\operatorname{prog}(v\in \Q_{\ge0},\ v+1 \le 0:\ \alpha v) = \top}
\end{stmt}
\end{lemma}

\begin{proof}
\begin{plines}
\pline{(1)}{The goal follows by the prices \(1\) of the rows.}
\end{plines}
\end{proof}

\begin{lemma}[{a program with an upper row}]\label{lem:progupper}
A row <= with a positive price: on v + w = 1 the cost 2 alpha - alpha v is least at v = 1/2, where the row v <= 1/2
is tight (price alpha); with the row = it leaves no direction, so the face is the vertex (1/2, 1/2).
\begin{stmt}
\sit{given}{\alpha\in \Q_{>0}}
\sit{(1)}{\operatorname{prog}(v,w\in \Q_{\ge0},\ v+w = 1,\ v \le 1/2:\ \alpha v+2 \alpha w) = (3 \alpha/2,1)}
\end{stmt}
\end{lemma}

\begin{proof}
\begin{plines}
\pline{(1)}{The goal follows by the prices \(2 \alpha, \alpha\) of the rows at the points \(v=1/2, w=1/2\).}
\end{plines}
\end{proof}

\begin{lemma}[{a program in a given}]\label{lem:proggiven}
The atom in a given: p is the pair of the program (its least value alpha at v = 1); no word reads the atom yet, so
the item checks only that the statement is formed.
\begin{stmt}
\sit{given}{\alpha\in \Q_{>0},\ p\in \mathrm{Pair},\ \operatorname{prog}(v\in \Q_{\ge0},\ v \ge 1:\ \alpha v) = p}
\sit{(1)}{\alpha/2 < \alpha}
\end{stmt}
\end{lemma}

\begin{proof}
\begin{plines}
\pline{(1)}{The goal follows from the facts at hand.}
\end{plines}
\end{proof}

\subsection{Coordinates and lattice points}\label{sec:lattice}

Two constructs of the list of primitives. The coordinate x_(i) of a tuple x of a
type row X^n (X one of NN, QQ>=0, QQ>0) is read under a binder over the coordinates: sum i, t; prod i, t;
forall i : P; exists i : P, where i runs over the coordinates of the tuples it indexes (Lean: x i with i : Fin n).
The index is only an index: no arithmetic on it, and no numeral as a coordinate. The sum sum u <= b, t runs over
the lattice points u of NN^n with u <= b coordinatewise, b a tuple of NN^n (Finset.Icc 0 b); rows cut the box
through iverson(...), so the number of points of a cut box is a sum of iverson(ROWS).

\begin{lemma}[{the points of a square under a diagonal}]\label{lem:latcount}
At fixed values the word count substitutes each tuple that a hypothesis fixes and evaluates the sums in the kernel.
\end{lemma}

\begin{lemma}[{the points of a square with a coordinate 1}]\label{lem:latsome}
A quantifier over the coordinates inside a row: the points of the square {0, 1}^2 off the corner 0.
\end{lemma}

\begin{lemma}[{the points of a box}]\label{lem:latcard}
The number of points of a box is the product of its sides plus one (Pi.card_Icc).
\end{lemma}

\begin{lemma}[{a cut box has at most the points of the box}]\label{lem:latbox}
The word box compares a count over a box cut by a row with the points of a box [0, c]: from above when the cut lies
under c (here c = b, the cut is a subset), from below when [0, c] lies in the cut.
\end{lemma}

\begin{lemma}[{a cut box under a corner}]\label{lem:latcorner}
From above by a smaller box: a point with coordinate sum at most 2 has every coordinate at most 2 (a term of a sum of
terms at least zero is at most the sum).
\end{lemma}

\begin{lemma}[{a box inside a cut box}]\label{lem:latunder}
From below: a box [0, c] under b lies in the cut by the content of c (each sum over the coordinates is monotone in
the point).
\end{lemma}

\begin{lemma}[{two sums over the coordinates}]\label{lem:latsum}
sum compares two sums over the coordinates term by term, with the hypotheses over every coordinate at the index.
\end{lemma}

\begin{lemma}[{a sum of coordinates at most 1}]\label{lem:latind}
ind over the coordinates: induction on the length n, the last coordinate split off (Fin.sum_univ_castSucc) and the
relation at the first n coordinates (Fin.init).
\end{lemma}

\begin{lemma}[{the points of a list of boxes}]\label{lem:latlist}
The coordinate of a member of a family of tuples, s_(k)_(i): a list of N boxes [0, s_(k)] has the points of each
box, the product of its sides plus one, summed over the list (the word box at each member, through sum).
\end{lemma}

\begin{lemma}[{two tuples of tuples entry by entry}]\label{lem:latrows}
A tuple of tuples (NN^n)^k, k rows of n coordinates: its entry
f_(iota)_(i), both indices bound over coordinates (Lean: f iota i with iota : Fin k, i : Fin n). Two tuples of
tuples compared entry by entry, by sum over the rows and sum over the coordinates.
\end{lemma}

\begin{lemma}[{a shift into a larger box}]\label{lem:latinj}
The word inj counts by an injective relation: each point of the cut box [0, b] is related to a point of the cut
box [0, c] (the totality, the relation the conjunction after exists), and two points related to one point are
equal (the injectivity), so the first cut has at most the points of the second (here the shift by 1 in every
coordinate). The exact model reads the quantifiers over NN^n on the boxes that their clauses forall i : u_(i) <= t
bound. The two hypotheses may come from with (with: the injectivity proved first).
\begin{stmt}
\sit{with}{n\in \N,\ b,c\in \mathrm{Grid}(n),\ (\forall u\in \mathrm{Grid}(n):\ (\forall i:\ u_i \le b_i),\ \textstyle\sum_i u_i \le 2\Rightarrow \exists v\in \mathrm{Grid}(n):\ (\forall i:\ v_i \le c_i),\ (\forall i:\ v_i = u_i+1),\ \textstyle\sum_i v_i \le 2+n):\ \textstyle\sum_{u\le b} \mathbf 1[\textstyle\sum_i u_i \le 2] \le \textstyle\sum_{v\le c} \mathbf 1[\textstyle\sum_i v_i \le 2+n]}
\end{stmt}
\end{lemma}

\begin{proof}
\begin{plines}
\pline{with}{\textup{(1)}~\(\forall u,w,v\in \mathrm{Grid}(n):\ (\forall i:\ u_i \le b_i),\ (\forall i:\ w_i \le b_i),\ (\forall i:\ v_i \le c_i),\ (\forall i:\ v_i = u_i+1),\ (\forall i:\ v_i = w_i+1)\Rightarrow \forall i:\ u_i = w_i\) by logic and from the facts at hand. The goal follows by an injective relation from the points of one count to those of the other.}
\end{plines}
\end{proof}

\begin{lemma}[{a content under its level, term by term}]\label{lem:latcellsum}
A sum over a cell: the points of a range at which every clause holds, each clause a comparison or a family of
comparisons. Over an index range, sum m in a..b : C_1, ..., C_J, t; over the coordinates, sum i : C_1, ..., C_J, t;
over the positive lattice points u of (NN>=1)^n, sum u : C_1, ..., C_J, t, its first clause the content against a
level, sum i, c_(i) * u_(i) + O rel L with every weight c_(i) at least 1, so that u_(i) <= L and the cell is finite
(the point may name its type, sum u in PosGrid(n) : ..., or sum u in Grid(n) : ... for the points of NN^n). A point
of a cell satisfies each clause there (h_ck_u, the clause k at the point u), and sum reads the terms at a point
with these facts.
\end{lemma}

\begin{lemma}[{a cell of coordinates, term by term}]\label{lem:latcellcoord}
Over the coordinates: the coordinates at which x lies below y.
\end{lemma}

\begin{lemma}[{a chamber and its wall}]\label{lem:latcellwall}
The word trichotomy: at each point a comparison A, B has exactly one of A < B, A = B, B < A, so a cell sum is the sum
over the three parts of one more clause, and a weak clause A <= B is its chamber A < B and its wall A = B.
\end{lemma}

\begin{lemma}[{a range cell in three parts}]\label{lem:latcellthree}
Over an index range: the members of a family below, at and above a count.
\end{lemma}

\begin{lemma}[{the points of a simplex wall}]\label{lem:latcellsimplex}
The word simplex on a cell: the points t of NN^(d + 1) with sum k, t_(k) = m are the weak compositions of m into
d + 1 parts, binom(m + d, d).
\end{lemma}

\begin{lemma}[{the positive points of a simplex wall}]\label{lem:latcellpossimplex}
Over the positive points: t in (NN>=1)^(d + 1) with sum k, t_(k) = m + d + 1 is s + 1 for a weak composition s of
m, so the wall at m + r has binom(m + d, d) points too.
\end{lemma}

\begin{lemma}[{a simplex wall inside its chamber}]\label{lem:latcellinj}
The word inj on cells: each point of the first cell is related to a point of the second (the totality, its premise
the clauses of the first cell, its last clauses those of the second), and two points related to one point are equal
(the injectivity, its premises the clauses at both points and the relation at both), so the first cell has at most
the points of the second. Each premise is a fact of a membership, and each image lies in the box of its level.
\end{lemma}

\begin{lemma}[{one cell read through a tie}]\label{lem:latcellsame}
The word same: two cells over one holder with one term have the same points when their clauses agree at each point
of the holder once the named ties are read there.
\end{lemma}

\begin{lemma}[{a range cell with one point}]\label{lem:latcellpoint}
A cell of one point: the clause m = p with a <= p <= b, the term at p.
\end{lemma}

\begin{lemma}[{a coordinate cell with one point}]\label{lem:latcellpointc}
Over the coordinates of x: the cell i = k, 1 <= k <= n, has the one point k.
\end{lemma}

\begin{lemma}[{a coordinate sum exchanged with a range cell}]\label{lem:latcellexch}
The exchange: a sum over the coordinates of a range cell is the range cell of the coordinate cells. Here the
meetings m of 1..n, from e_(m) to f_(m), and P_(i)_(j) the meetings from i to j.
\end{lemma}

\begin{lemma}[{a part left out by the range}]\label{lem:latcellbound}
A part left out of trichotomy is empty: above the range (n < m against m <= n), or refuted at each point by a
certificate of 1 <= 0 from the clauses there, the families over the range at the point and the comparison.
\end{lemma}

\begin{lemma}[{parts left out by a certificate}]\label{lem:latcellcert}
\end{lemma}

\begin{lemma}[{a sum split into the cells of one comparison}]\label{lem:latcellplain}
A sum with no clause is the cell of no clause on its holder, so trichotomy splits it too: over the coordinates at an
index j, below j, at j and above j; over a range at its top; and by a certificate when the family has one value.
\begin{stmt}
\sit{range}{n\in \N_{\ge1},\ e_1,\dots,e_n\in \N_{\ge1}:\ \textstyle\sum_{m=1}^n e_m = \textstyle\sum_{\substack{1\le m\le n\\ m < n}} e_m+(\textstyle\sum_{\substack{1\le m\le n\\ m = n}} e_m)}
\sit{cert}{n\in \N_{\ge1},\ e_1,\dots,e_n\in \N_{\ge1},\ (\forall m\in \N_{\ge1}:\ m \le n\Rightarrow e_m = 1):\ \textstyle\sum_{m=1}^n e_m = \textstyle\sum_{\substack{1\le m\le n\\ e_m = 1}} e_m}
\end{stmt}
\end{lemma}

\begin{proof}
\begin{plines}
\pline{range}{The goal follows splitting the cell at each point by the three cases of one comparison.}
\pline{cert}{Splitting the cell at each point by the three cases of one comparison, each part left out is empty: Clause 1: the goal follows combining \(h_{tri_m}\), with \((\mathrm{H}_{1.m})\) read from right to left. Clause 2: the goal follows combining \(h_{tri_m}\) and \((\mathrm{H}_{1.m})\).}
\end{plines}
\end{proof}

\begin{lemma}[{a coordinate cell at an index}]\label{lem:latcellpointidx}
A coordinate cell whose one point is an index j: the term at j.
\end{lemma}

\begin{lemma}[{a cell read up to the order of its clauses}]\label{lem:latcellorder}
A cell is the conjunction of its clauses, so two cells with the same clauses in another order are one cell (same,
with no tie); read through ties, the clauses agree up to their order.
\begin{stmt}
\sit{tie}{n,i,j\in \N_{\ge1},\ e_1,\dots,e_n,f_1,\dots,f_n\in \N_{\ge1},\ (\forall m\in \N_{\ge1}:\ m \le n\Rightarrow e_m = m):\ \textstyle\sum_{\substack{1\le m\le n\\ e_m = j\\ f_m = i}} 1 = \textstyle\sum_{\substack{1\le m\le n\\ f_m = i\\ m = j}} 1}
\end{stmt}
\end{lemma}

\begin{proof}
\begin{plines}
\pline{tie}{The goal follows reading the clauses of the two cells as one by \((\mathrm{H}_{1})\).}
\end{plines}
\end{proof}

\begin{lemma}[{a point of a range cell lies in its range}]\label{lem:latcellrange}
A point of a range cell is a point of the range, so the families over the range hold there: here e_(m) <= n, so
that the coordinate cell e_(m) = j has its one point.
\end{lemma}

\begin{lemma}[{a range cell with one point among its clauses}]\label{lem:latcellpointmany}
A cell of one point among several clauses: the clause m = p, the other clauses true at p, the term at p.
\begin{stmt}
\sit{coord}{n\in \N_{\ge1},\ x\in \mathrm{Point}(n),\ (\forall i:\ x_i \le 1):\ \forall i:\ \textstyle\sum_{\substack{1\le m\le n\\ 1 \le m\\ m = i}} 1 = 1}
\end{stmt}
\end{lemma}

\begin{proof}
\begin{plines}
\pline{coord}{The goal follows by logic and term by term.}
\end{plines}
\end{proof}

\begin{lemma}[{a triangle summed in either order}]\label{lem:latcelltri}
The exchange of two sums over ranges through a cell that relates both indices: the pairs l <= i of a triangle
counted by rows or by columns.
\begin{stmt}
\sit{gap}{r\in \N_{\ge1},\ e_1,\dots,e_r\in \N_{\ge1}:\ \textstyle\sum_{h=0}^r (\textstyle\sum_{\substack{1\le g\le r\\ h < g}} e_g) = \textstyle\sum_{g=1}^r (\textstyle\sum_{\substack{0\le h\le r\\ h < g}} e_g)}
\end{stmt}
\end{lemma}

\begin{proof}
\begin{plines}
\pline{gap}{The goal follows term by term.}
\end{plines}
\end{proof}

\begin{lemma}[{the bounds of a power difference}]\label{lem:zpowdiff}
The word powdiff bounds a power difference: from the tie l^p + p * W = (l + s)^p with a <= l and l + s <= b, W lies
between a^(p - 1) * s and b^(p - 1) * s, written without subtraction (times a, and times b); the geometric identity
(sum_(i < p) (s + l)^i * l^(p - 1 - i)) * s + l^p = (s + l)^p, each of its p terms between a^(p - 1) and b^(p - 1).
Under binders over the rows and the coordinates of tuples of tuples, the hypotheses over the same coordinates are
read at the indices of the goal.
\end{lemma}

\begin{lemma}[{the bounds of power differences entry by entry}]\label{lem:latpowdiff}
\end{lemma}

\begin{lemma}[{weak duality from a tie of forms}]\label{lem:latdual}
The word dual is weak duality from a tie of forms: the constraints P_(iota) . v >= N_(iota) . v + e_(iota), each
times its price y_(iota) >= 0 and summed, with the two sums swapped, give c . v >= y . e by the tie at each
coordinate, c_(i) + sum_iota y_(iota) * N_(iota)_(i) = sum_iota y_(iota) * P_(iota)_(i); the forms are on both
sides, so no subtraction enters. With an offset a_(iota) on the larger side of each constraint, the goal carries
y . a on the left (offset); the constraints are the universal whose two sides read a sum over the coordinates
(pick: another universal comes first), the first whose smaller side adds the e_(iota) of the goal (match, among the
parts of a conjunction), natural constraints read in QQ through their casts (nat). With a form d . v on the right, the tie reads c_(i) + sum_iota y_(iota) * N_(iota)_(i)
= d_(i) + sum_iota y_(iota) * P_(iota)_(i) and the goal carries d . v on the right (form).
\begin{stmt}
\sit{offset}{k,n\in \N,\ P,N\in \mathrm{Forms}(k,n),\ y\in \mathrm{Point}(k),\ e,a\in \mathrm{Grid}(k),\ c\in \mathrm{Grid}(n),\ (\forall i:\ c_i+(\textstyle\sum_{\iota} y_{\iota} N_{\iota,i}) = \textstyle\sum_{\iota} y_{\iota} P_{\iota,i}):\ \forall v\in \mathrm{Point}(n):\ (\forall \iota:\ \textstyle\sum_i P_{\iota,i} v_i+a_{\iota} \ge \textstyle\sum_i N_{\iota,i} v_i+e_{\iota})\Rightarrow \textstyle\sum_i c_i v_i+(\textstyle\sum_{\iota} y_{\iota} a_{\iota}) \ge \textstyle\sum_{\iota} y_{\iota} e_{\iota}}
\sit{form}{k,n\in \N,\ P,N\in \mathrm{Forms}(k,n),\ y\in \mathrm{Point}(k),\ e,a\in \mathrm{Grid}(k),\ c,d\in \mathrm{Grid}(n),\ (\forall i:\ c_i+(\textstyle\sum_{\iota} y_{\iota} N_{\iota,i}) = d_i+(\textstyle\sum_{\iota} y_{\iota} P_{\iota,i})):\ \forall v\in \mathrm{Point}(n):\ (\forall \iota:\ \textstyle\sum_i P_{\iota,i} v_i+a_{\iota} \ge \textstyle\sum_i N_{\iota,i} v_i+e_{\iota})\Rightarrow \textstyle\sum_i c_i v_i+(\textstyle\sum_{\iota} y_{\iota} a_{\iota}) \ge \textstyle\sum_i d_i v_i+(\textstyle\sum_{\iota} y_{\iota} e_{\iota})}
\sit{pick}{k,n\in \N,\ P,N\in \mathrm{Forms}(k,n),\ y\in \mathrm{Point}(k),\ e\in \mathrm{Grid}(k),\ c\in \mathrm{Grid}(n),\ (\forall \iota:\ y_{\iota} \le y_{\iota}+e_{\iota}),\ (\forall i:\ c_i+(\textstyle\sum_{\iota} y_{\iota} N_{\iota,i}) = \textstyle\sum_{\iota} y_{\iota} P_{\iota,i}):\ \forall v\in \mathrm{Point}(n):\ (\forall \iota:\ \textstyle\sum_i P_{\iota,i} v_i \ge \textstyle\sum_i N_{\iota,i} v_i+e_{\iota})\Rightarrow \textstyle\sum_i c_i v_i \ge \textstyle\sum_{\iota} y_{\iota} e_{\iota}}
\sit{nat}{k,n\in \N,\ P,N\in \mathrm{Rows}(k,n),\ y\in \mathrm{Point}(k),\ e\in \mathrm{Grid}(k),\ c\in \mathrm{Grid}(n),\ (\forall i:\ c_i+(\textstyle\sum_{\iota} y_{\iota} N_{\iota,i}) = \textstyle\sum_{\iota} y_{\iota} P_{\iota,i}):\ \forall v\in \mathrm{Grid}(n):\ (\forall \iota:\ \textstyle\sum_i P_{\iota,i} v_i \ge \textstyle\sum_i N_{\iota,i} v_i+e_{\iota})\Rightarrow \textstyle\sum_i c_i v_i \ge \textstyle\sum_{\iota} y_{\iota} e_{\iota}}
\sit{match}{k,n\in \N,\ P,N\in \mathrm{Rows}(k,n),\ y\in \mathrm{Point}(k),\ e,f\in \mathrm{Grid}(k),\ c\in \mathrm{Grid}(n),\ (\forall i:\ c_i+(\textstyle\sum_{\iota} y_{\iota} N_{\iota,i}) = \textstyle\sum_{\iota} y_{\iota} P_{\iota,i}):\ \forall v\in \mathrm{Grid}(n):\ (\forall \iota:\ \textstyle\sum_i P_{\iota,i} v_i \ge \textstyle\sum_i N_{\iota,i} v_i+f_{\iota}),\ (\forall \iota:\ \textstyle\sum_i P_{\iota,i} v_i \ge \textstyle\sum_i N_{\iota,i} v_i+e_{\iota})\Rightarrow \textstyle\sum_i c_i v_i \ge \textstyle\sum_{\iota} y_{\iota} e_{\iota}}
\end{stmt}
\end{lemma}

\begin{proof}
\begin{plines}
\pline{offset}{The goal follows by weak duality.}
\pline{form}{The goal follows by weak duality.}
\pline{pick}{The goal follows by weak duality.}
\pline{nat}{The goal follows by weak duality.}
\pline{match}{The goal follows by weak duality.}
\end{plines}
\end{proof}

\begin{stmt}
\sit{relation}{l\ \text{zrows}\ t:\iff k,n\in \N,\ l\in \mathrm{Forms}(k,n),\ t\in \Q_{\ge0},\ \forall \iota:\ \forall i:\ l_{\iota,i} \le t}
\end{stmt}

\begin{lemma}[{witnesses of tuples of tuples}]\label{lem:zwit}
A relation in a conclusion is proved by its fields, one word each in the order of the row (or one per row, the
bounds of its declarations closed by arithmetic). The word take gives the witnesses of the existentials that lead a
goal, in their order: a number by a term, a tuple of tuples by a formula in its indices (each coordinate with its
sign) or by join(a, b), the rows of a then those of b; the word after the semicolon proves the rest. A goal over the
rows of a join is split at the seam: a forall into its two halves, each by the same word, a sum into the sum of
the two sums, an exists into the disjunction of the two.
\begin{stmt}
\sit{field}{k,n\in \N,\ l\in \mathrm{Forms}(k,n),\ t\in \Q_{\ge0},\ l\ \text{zrows}\ t:\ l\ \text{zrows}\ t+1}
\sit{scale}{k,n\in \N,\ l\in \mathrm{Forms}(k,n),\ t\in \Q_{\ge0},\ h\in \Q_{>0},\ l\ \text{zrows}\ t:\ \exists m\in \mathrm{Forms}(k,n):\ m\ \text{zrows}\ h t}
\sit{join}{k,k',n\in \N,\ l\in \mathrm{Forms}(k,n),\ l'\in \mathrm{Forms}(k',n),\ t\in \Q_{\ge0},\ l\ \text{zrows}\ t,\ l'\ \text{zrows}\ t:\ \exists r\in \N:\ \exists m\in \mathrm{Forms}(r,n):\ m\ \text{zrows}\ t,\ \textstyle\sum_{\iota} (\textstyle\sum_i m_{\iota,i}) = \textstyle\sum_{\iota} (\textstyle\sum_i l_{\iota,i})+(\textstyle\sum_{\iota} (\textstyle\sum_i l'_{\iota,i}))}
\end{stmt}
\end{lemma}

\begin{proof}
\begin{plines}
\pline{field}{The goal follows from the facts at hand.}
\pline{scale}{Take \(m_{\iota,i}=h l_{\iota,i}\). The goal follows from the facts at hand.}
\pline{join}{Take \(r=k+k', m=\mathrm{join}(l,l')\). The parts of the goal follow: (i) from the facts at hand; (ii) by expanding.}
\end{plines}
\end{proof}

\begin{lemma}[{first-order steps}]\label{lem:zlogic}
The word logic proves a goal that first-order logic gives from the hypotheses: each hypothesis read as a rule
(binders over points, objects and indices, memberships, premises) is applied when its premises are facts, the
first matched with a fact the goal or an earlier step brought in (a premise, the membership of a point), the others
found among the facts; its conclusion is destructed (an existential gives its witnesses, a conjunction its parts)
or kept as a fact; a hypothesis S subseteq T is applied to each membership in S (a membership of a union read as a
disjunction); up to three rounds, then a tautology. Here a point of C has a point of B under it, so one of B'; and a
cover of B by two regions gives a cover by the join of their two tuples of tuples, at the seam the disjunction of
the two existentials.
\begin{stmt}
\sit{inside}{n\in \N,\ B,B',C\in \mathrm{Region}(n),\ B \subseteq B',\ (\forall w\in C:\ \exists v\in B:\ \forall i:\ v_i \le w_i):\ \forall w\in C:\ \exists v\in B':\ \forall i:\ v_i \le w_i}
\sit{cover}{k,k',n\in \N,\ l\in \mathrm{Forms}(k,n),\ l'\in \mathrm{Forms}(k',n),\ B,B_1,B_2\in \mathrm{Region}(n),\ B \subseteq B_1\cup B_2,\ (\forall w\in B_1:\ \exists \iota:\ \forall i:\ l_{\iota,i} \le w_i),\ (\forall w\in B_2:\ \exists \iota:\ \forall i:\ l'_{\iota,i} \le w_i):\ \exists r\in \N:\ \exists m\in \mathrm{Forms}(r,n):\ \forall w\in B:\ \exists \iota:\ \forall i:\ m_{\iota,i} \le w_i}
\end{stmt}
\end{lemma}

\begin{proof}
\begin{plines}
\pline{inside}{The goal follows by logic.}
\pline{cover}{Take \(r=k+k', m=\mathrm{join}(l,l')\). The goal follows by logic.}
\end{plines}
\end{proof}

\begin{lemma}[{a power under the sum of the powers}]\label{lem:zindpow}
ind with the signs of its sums: the step from n to n + 1 adds the term r^(n + 1), at least 0 from 0 <= r
(pow_nonneg), and the sum at n is at least 0 (Finset.sum_nonneg).
\end{lemma}

\begin{lemma}[{forward steps}]\label{lem:zstep}
logic[W] proves the rest by the word W, which reads the witnesses by the names their existentials bind (with a
suffix _2, _3, ... when a name is taken); a universal over a type with no premise is instantiated at each object of
that type in scope. with[CLAUSE by W; W'] is a step: the clause, in the names in scope after the binders and
premises of the goal are introduced, proved by W, then a fact for W' (an existential destructed). take gives an
index, or a point of a region (with its membership). Here: the level of a dilation, 2^D * K'(w') = K(w) with
2^L * K(w) <= 1, the powers split and the equation times 2^L; a premise derived first by with; and a preimage
u_(i) = w_(i) / d_(i), a point of B at u, then one of B' at d * u = w.
\begin{stmt}
\sit{level}{n,L,D\in \N,\ d\in \mathrm{Weight}(n),\ K,K'\in \mathrm{Loss}(n),\ B,B'\in \mathrm{Region}(n),\ (\forall w\in \mathrm{Point}(n):\ \forall w'\in \mathrm{Point}(n):\ (\forall i:\ w'_i = d_i w_i)\Rightarrow 2^D K'(w') = K(w)),\ (\forall w'\in B':\ \exists w\in B:\ \forall i:\ w'_i = d_i w_i),\ (\forall w\in B:\ 2^L K(w) \le 1):\ \forall w'\in B':\ 2^{L+D} K'(w') \le 1}
\sit{premise}{k,n,L\in \N,\ l\in \mathrm{Forms}(k,n),\ t\in \Q_{\ge0},\ K\in \mathrm{Loss}(n),\ B\in \mathrm{Region}(n),\ t\cdot 2^L \le 1,\ (\forall w\in B:\ K(w) \le t),\ (\forall w\in B:\ 2^L K(w) \le 1\Rightarrow \exists \iota:\ \forall i:\ l_{\iota,i} \le w_i):\ \forall w\in B:\ \exists \iota:\ \forall i:\ l_{\iota,i} \le w_i}
\sit{preimage}{n\in \N,\ d\in \mathrm{Weight}(n),\ B,B'\in \mathrm{Region}(n),\ (\forall w\in \mathrm{Point}(n):\ \exists v\in B:\ \forall i:\ v_i = w_i),\ (\forall v\in B:\ \exists v'\in B':\ \forall i:\ v'_i = d_i v_i):\ \forall w\in \mathrm{Point}(n):\ \exists y\in B':\ \forall i:\ y_i = w_i}
\end{stmt}
\end{lemma}

\begin{proof}
\begin{plines}
\pline{level}{The goal follows by logic and from the facts at hand.}
\pline{premise}{By logic: \textup{(1)}~\(2^L K(w) \le 1\) from the facts at hand. By logic: take \(\iota=\iota\). The goal follows from the facts at hand.}
\pline{preimage}{\textup{(1)}~\(\exists u\in \mathrm{Point}(n):\ \forall i:\ w_i = d_i u_i\). Take \(u_i=w_i/d_i\). The goal follows from the facts at hand. This proves \textup{(1)}. By logic: take \(y=v'\). The goal follows from the facts at hand.}
\end{plines}
\end{proof}

\begin{lemma}[{a factor into a sum, two products merged, brackets by their conditions}]\label{lem:zsumfac}
sum[W] first brings a factor free of the index into its sum (c * sum = sum of c times the terms) and merges two
products over the same coordinates (prod A times prod B = prod of A * B), then compares term by term, products as
sums; an equation of two Iverson brackets is the equivalence of their conditions, one word each way. Two products
are compared term by term, = or <= (each factor of the smaller product at least 0). Two nested sums in the other
order on the right are swapped first (over ranges, coordinates or the lattice points of a box). Two sums over the
same range added are one sum of the two terms, on either side (addcoords, addrange, addle). Over a box, the facts
over the points of its type are read at the index (boxfact, boxwith). A product of powers of one base is one power
of the sum of the exponents (powsum), then W proves the relation of the powers (powd, powle).
\begin{stmt}
\sit{factor}{k,n\in \N,\ l\in \mathrm{Forms}(k,n),\ c\in \mathrm{Point}(n):\ \textstyle\sum_{\iota} (\textstyle\prod_i c_i l_{\iota,i}) = (\textstyle\prod_i c_i)\cdot (\textstyle\sum_{\iota} (\textstyle\prod_i l_{\iota,i}))}
\sit{bracket}{k,n\in \N,\ l\in \mathrm{Forms}(k,n),\ t\in \Q_{\ge0}:\ \textstyle\sum_{\iota} \mathbf 1[\forall i:\ 2 l_{\iota,i} < 2 t] = \textstyle\sum_{\iota} \mathbf 1[\forall i:\ l_{\iota,i} < t]}
\sit{prodle}{k,n\in \N,\ V,W\in \mathrm{Forms}(k,n),\ A\in \mathrm{Point}(n),\ (\forall \iota:\ \forall i:\ V_{\iota,i} \le A_i W_{\iota,i}):\ \textstyle\sum_{\iota} (\textstyle\prod_i V_{\iota,i}) \le \textstyle\sum_{\iota} (\textstyle\prod_i A_i W_{\iota,i})}
\sit{prodge}{n\in \N,\ a,b\in \mathrm{Point}(n),\ (\forall i:\ a_i \le b_i):\ \textstyle\prod_i (b_i+1) \ge \textstyle\prod_i a_i}
\sit{swapbox}{n,J\in \N,\ b\in \mathrm{Grid}(n):\ \textstyle\sum_{j=0}^J (\textstyle\sum_{u\le b} \mathbf 1[\textstyle\sum_i u_i \le j]) = \textstyle\sum_{u\le b} (\textstyle\sum_{j=0}^J \mathbf 1[\textstyle\sum_i u_i \le j])}
\sit{swapcoords}{k,n\in \N,\ l\in \mathrm{Forms}(k,n):\ \textstyle\sum_{\iota} (\textstyle\sum_i l_{\iota,i}) = \textstyle\sum_i (\textstyle\sum_{\iota} l_{\iota,i})}
\sit{swaple}{n,J\in \N,\ a\in \mathrm{Point}(n):\ \textstyle\sum_{j=0}^J (\textstyle\sum_i a_i) \le \textstyle\sum_i (\textstyle\sum_{j=0}^J (a_i+1))}
\sit{swapweight}{n,J\in \N,\ b\in \mathrm{Grid}(n),\ w_0,\dots,w_J\in \Q_{\ge0}:\ \textstyle\sum_{u\le b} (\textstyle\sum_{j=0}^J \mathbf 1[\textstyle\sum_i u_i \le j] w_j) = \textstyle\sum_{j=0}^J (\textstyle\sum_{u\le b} \mathbf 1[\textstyle\sum_i u_i \le j]) w_j}
\sit{addcoords}{n\in \N,\ \mu,F,o\in \mathrm{Point}(n):\ \textstyle\sum_{\iota} \mu_{\iota}\cdot (F_{\iota}+o_{\iota}) = \textstyle\sum_{\iota} \mu_{\iota} F_{\iota}+(\textstyle\sum_{\kappa} \mu_{\kappa} o_{\kappa})}
\sit{addrange}{J\in \N:\ \textstyle\sum_{j=0}^J (j+1) = \textstyle\sum_{j=0}^J j+(\textstyle\sum_{j=0}^J 1)}
\sit{addle}{n\in \N,\ \mu,F,o\in \mathrm{Point}(n):\ \textstyle\sum_{\iota} \mu_{\iota} F_{\iota}+(\textstyle\sum_{\iota} \mu_{\iota} o_{\iota}) \le \textstyle\sum_{\iota} \mu_{\iota}\cdot (F_{\iota}+o_{\iota}+1)}
\sit{boxfact}{n\in \N,\ b\in \mathrm{Grid}(n),\ (\forall u\in \mathrm{Grid}(n):\ \mathbf 1[\textstyle\sum_i u_i \le 2] \le \mathbf 1[\textstyle\sum_i u_i \le 3]):\ \textstyle\sum_{u\le b} \mathbf 1[\textstyle\sum_i u_i \le 2] \le \textstyle\sum_{u\le b} \mathbf 1[\textstyle\sum_i u_i \le 3]}
\sit{boxwith}{n\in \N,\ b\in \mathrm{Grid}(n),\ (\forall u\in \mathrm{Grid}(n):\ \mathbf 1[\textstyle\sum_i u_i \le 2] \le \mathbf 1[\textstyle\sum_i u_i \le 3]):\ \textstyle\sum_{w\le b} \mathbf 1[\textstyle\sum_i w_i \le 2] \le \textstyle\sum_{w\le b} \mathbf 1[\textstyle\sum_i w_i \le 3]}
\sit{powsum}{k\in \N,\ x\in \Q_{>0},\ a\in \mathrm{Grid}(k):\ \textstyle\prod_{\iota} x^{a_{\iota}} = x^{\textstyle\sum_{\iota} a_{\iota}}}
\sit{powd}{k,d\in \N,\ x\in \Q_{>0},\ a\in \mathrm{Grid}(k),\ d = \textstyle\sum_{\iota} a_{\iota}:\ \textstyle\prod_{\iota} x^{a_{\iota}} = x^d}
\sit{powle}{k,d\in \N,\ a\in \mathrm{Grid}(k),\ \textstyle\sum_{\iota} a_{\iota} \le d:\ \textstyle\prod_{\iota} 2^{a_{\iota}} \le 2^d}
\end{stmt}
\end{lemma}

\begin{proof}
\begin{plines}
\pline{factor}{The goal follows term by term, term by term.}
\pline{bracket}{The goal follows term by term, from the facts at hand (each part).}
\pline{prodle}{The goal follows term by term, term by term, from the facts at hand.}
\pline{prodge}{The goal follows term by term, from the facts at hand.}
\pline{swapbox}{The goal follows term by term.}
\pline{swapcoords}{The goal follows term by term.}
\pline{swaple}{The goal follows term by term, term by term, from the facts at hand.}
\pline{swapweight}{The goal follows term by term, term by term.}
\pline{addcoords}{The goal follows term by term.}
\pline{addrange}{The goal follows term by term.}
\pline{addle}{The goal follows term by term.}
\pline{boxfact}{The goal follows term by term, from the facts at hand.}
\pline{boxwith}{\textup{(1)}~\(\forall v\in \mathrm{Grid}(n):\ \mathbf 1[\textstyle\sum_i v_i \le 2] \le \mathbf 1[\textstyle\sum_i v_i \le 3]\) by logic. The goal follows term by term, from the facts at hand.}
\pline{powsum}{The goal follows term by term.}
\pline{powd}{The goal follows term by term, from the facts at hand.}
\pline{powle}{The goal follows term by term, from the facts at hand.}
\end{plines}
\end{proof}

\begin{lemma}[{a term under its sum}]\label{lem:zterm}
term[J] bounds a sum from below by its term at the index J: each term is at least 0, its sign built from the signs
of its atoms, and J lies in the range of the sum; term[J; W] proves first that t is at most the term at J. Over
the coordinates of tuples, J is a coordinate index in scope (the witness of an existential: coord, coordsame).
\begin{stmt}
\sit{below}{k,J\in \N:\ (J+1)^k/2^{J+1} \le \textstyle\sum_{j=0}^J (j+1)^k/2^j}
\sit{coord}{k\in \N,\ p,s\in \mathrm{Point}(k),\ (\exists j:\ 2 \le p_j s_j):\ 2 \le \textstyle\sum_{\iota} p_{\iota} s_{\iota}}
\sit{coordsame}{k\in \N,\ p,s\in \mathrm{Point}(k),\ (\exists j:\ 2 \le p_j s_j):\ \exists j:\ p_j s_j \le \textstyle\sum_{\iota} p_{\iota} s_{\iota}}
\end{stmt}
\end{lemma}

\begin{proof}
\begin{plines}
\pline{below}{The goal follows from the term at \(J\), from the facts at hand.}
\pline{coord}{\textup{(1)}~\(\exists j:\ 2 \le p_j s_j\) by logic. The goal follows from the term at \(j\).}
\pline{coordsame}{\textup{(1)}~\(\exists j:\ 2 \le p_j s_j\) by logic. Take \(j=j\). The goal follows from the term at \(j\).}
\end{plines}
\end{proof}

\begin{lemma}[{a telescope of ties, a count cancelled}]\label{lem:zcells}
ind with a hypothesis over every natural number (the tie of each term): the first case reads it at the least
index, the step at the new index, each case an exact combination with the relation at n. A count that multiplies
is at least 1, so it cancels: x + q * W = y = x + q * Z gives W = Z.
\begin{stmt}
\sit{telescope}{q\in \N_{\ge1},\ a,t\in \N,\ c_0,\dots,c_t\in \Q_{\ge0},\ (\forall j\in \N:\ (a+j)^q+q c_j = (a+j+1)^q):\ a^q+q\cdot (\textstyle\sum_{j=0}^t c_j) = (a+t+1)^q}
\sit{cancel}{q\in \N_{\ge1},\ W,Z,x,y\in \Q_{\ge0},\ x+q W = y,\ x+q Z = y:\ W = Z}
\end{stmt}
\end{lemma}

\begin{proof}
\begin{plines}
\pline{telescope}{The goal follows by induction on the range of the sums.}
\pline{cancel}{The goal follows from the facts at hand.}
\end{plines}
\end{proof}

cells compares two lists of boxes by their weights (each side weight W given by the tie l^q + q W = (l + s)^q):
the first list separated (each point in the interior of at most one box) and covered by the second (each point of
a box of the first in a box of the second), counted on the cells of a common grid; the hypotheses are found by
their shape, here one tie in a relation and one as it is.

\begin{stmt}
\sit{relation}{W\ \text{zweight}\ q\ \text{zcorners}\ l\ \text{zsides}\ s:\iff k,n\in \N,\ W,l\in \mathrm{Forms}(k,n),\ s\in \mathrm{Sides}(k,n),\ q\in \mathrm{Grid}(n),\ \forall \iota:\ \forall i:\ l_{\iota,i}^{q_i}+q_i W_{\iota,i} = (l_{\iota,i}+s_{\iota,i})^{q_i}}
\end{stmt}

\begin{lemma}[{two lists of boxes compared on the cells of a grid}]\label{lem:zcellsum}
\begin{stmt}
\sit{lists}{N,N',n\in \N,\ q\in \mathrm{Grid}(n),\ l,W\in \mathrm{Forms}(N,n),\ s\in \mathrm{Sides}(N,n),\ l',W'\in \mathrm{Forms}(N',n),\ s'\in \mathrm{Sides}(N',n),\ (\forall i:\ 1 \le q_i),\ W\ \text{zweight}\ q\ \text{zcorners}\ l\ \text{zsides}\ s,\ (\forall j:\ \forall i:\ l'_{ji}^{q_i}+q_i W'_{ji} = (l'_{ji}+s'_{ji})^{q_i}),\ (\forall w\in \mathrm{Point}(n):\ \textstyle\sum_k \mathbf 1[\forall i:\ l_{ki} < w_i < l_{ki}+s_{ki}] \le 1),\ (\forall k:\ \forall w\in \mathrm{Point}(n):\ (\forall i:\ l_{ki} \le w_i \le l_{ki}+s_{ki})\Rightarrow \exists j:\ \forall i:\ l'_{ji} \le w_i \le l'_{ji}+s'_{ji}):\ \textstyle\sum_k (\textstyle\prod_i W_{ki}) \le \textstyle\sum_k (\textstyle\prod_i W'_{ki})}
\end{stmt}
\end{lemma}

\begin{proof}
\begin{plines}
\pline{lists}{The goal follows by the weights of the boxes.}
\end{plines}
\end{proof}

\begin{lemma}[{induction on a natural binder}]\label{lem:zindn}
ind[n; W; W'] is an induction on the natural binder n for any goal: W proves the first case, W' the step, which
reads the case at n as a fact (an existential destructed, its witness under the name it binds); a binder of NN>=k
starts at k. Here: n + 1 <= 2^n; an existential whose witness grows by one at each step; 2 n <= 2^n from n = 2.
ind with no binder, on a product over the coordinates, splits off the last factor (a coordinate in an exponent
indexes its tuple).
\begin{stmt}
\sit{pow}{n\in \N:\ n+1 \le 2^n}
\sit{exist}{n\in \N:\ \exists m\in \N:\ n \le 2^m}
\sit{count}{n\in \N_{\ge2}:\ 2 n \le 2^n}
\sit{dil}{n,l\in \N,\ A\in \mathrm{Grid}(n):\ \textstyle\prod_i 2^{l A_i} = 2^{l\cdot (\textstyle\sum_i A_i)}}
\end{stmt}
\end{lemma}

\begin{proof}
\begin{plines}
\pline{pow}{By induction on \(n\). Base: the goal follows from the facts at hand. Step: the goal follows from the facts at hand.}
\pline{exist}{By induction on \(n\). Base: take \(m=0\). The goal follows from the facts at hand. Step: by logic: take \(m=m+1\). The goal follows from the facts at hand.}
\pline{count}{By induction on \(n\). Base: the goal follows from the facts at hand. Step: the goal follows from the facts at hand.}
\pline{dil}{The goal follows by induction on the range of the sums.}
\end{plines}
\end{proof}

\begin{lemma}[{a natural bound from a rational one}]\label{lem:zround}
A natural bound from a rational one: k <= x < j + 1 gives k <= j (the cast, then k < j + 1), in the closer and as
the side of mono (2^k <= 2^j); k = x = j gives k = j (the cast of the equation). A natural relation with a sum is
read in QQ with its casts at the atoms (sum), a sum or a product of its index too (index, prod).
\begin{stmt}
\sit{round}{k,j\in \N,\ x\in \Q_{\ge0},\ k \le x,\ x < j+1:\ k \le j}
\sit{strict}{k,j\in \N,\ x\in \Q_{\ge0},\ k < x,\ x \le j:\ k < j}
\sit{power}{k,j\in \N,\ x\in \Q_{\ge0},\ k \le x,\ x < j+1:\ 2^k \le 2^j}
\sit{shift}{\kappa,\kappa',c\in \N,\ x\in \Q_{\ge0},\ \kappa < x+1,\ x \le \kappa'+c:\ \kappa \le \kappa'+c}
\sit{eq}{k,j\in \N,\ x\in \Q_{\ge0},\ k = x,\ j = x:\ k = j}
\sit{sum}{n,c\in \N,\ a\in \mathrm{Grid}(n),\ x\in \Q_{\ge0},\ \textstyle\sum_i a_i \le c,\ x \le \textstyle\sum_i a_i:\ x \le c}
\sit{index}{J,c\in \N,\ x\in \Q_{\ge0},\ \textstyle\sum_{j=0}^J j \le c,\ x \le \textstyle\sum_{j=0}^J j:\ x \le c}
\sit{prod}{m,d\in \N,\ x\in \Q_{\ge0},\ (m+1)^d \le (\textstyle\prod_{i=1}^d i) \tbinom{m+d}{d},\ x \le (m+1)^d:\ x \le (\textstyle\prod_{i=1}^d i) \tbinom{m+d}{d}}
\end{stmt}
\end{lemma}

\begin{proof}
\begin{plines}
\pline{round}{The goal follows from the facts at hand.}
\pline{strict}{The goal follows from the facts at hand.}
\pline{power}{The goal follows by monotonicity.}
\pline{shift}{The goal follows from the facts at hand.}
\pline{eq}{The goal follows from the facts at hand.}
\pline{sum}{The goal follows from the facts at hand.}
\pline{index}{The goal follows from the facts at hand.}
\pline{prod}{The goal follows from the facts at hand.}
\end{plines}
\end{proof}

\begin{lemma}[{the deficit of two natural numbers}]\label{lem:zdeficit}
The deficit max(a - b, 0) of two natural numbers is a natural number (over N, the truncated difference a - b):
omega closes a natural goal with it (nat), it gives a tuple of NN by a formula (slack), and read in QQ it is the
same number (rat).
\begin{stmt}
\sit{nat}{a,b\in \N:\ \max(a-b,0)+b \ge a}
\sit{slack}{n\in \N,\ a,b\in \mathrm{Grid}(n):\ \exists s\in \mathrm{Grid}(n):\ \forall i:\ s_i+b_i \ge a_i}
\sit{rat}{a,b\in \N,\ x\in \Q_{\ge0},\ x = \max(a-b,0):\ x+b \ge a}
\end{stmt}
\end{lemma}

\begin{proof}
\begin{plines}
\pline{nat}{The goal follows from the facts at hand.}
\pline{slack}{Take \(s_i=\max(a_i-b_i,0)\). The goal follows from the facts at hand.}
\pline{rat}{The goal follows from the facts at hand.}
\end{plines}
\end{proof}

\begin{lemma}[{a reference at a natural sum}]\label{lem:zcastref}
A reference at a compound natural term (a sum) carries its cast whole, where the goal has the casts at the atoms:
the instance and its hypotheses are read up to casts (norm_cast).
\begin{stmt}
\sit{base}{m\in \N,\ x\in \Q_{\ge0},\ x \le m:\ x \le m+1}
\sit{sum}{n\in \N,\ a\in \mathrm{Grid}(n),\ x\in \Q_{\ge0},\ x \le \textstyle\sum_i a_i:\ x \le \textstyle\sum_i a_i+1}
\end{stmt}
\end{lemma}

\begin{proof}
\begin{plines}
\pline{base}{The goal follows from the facts at hand.}
\pline{sum}{The goal follows by \cref{lem:zcastref}(base).}
\end{plines}
\end{proof}

\begin{lemma}[{eq after a substitution}]\label{lem:zeqsub}
After the substitutions of eq (each defining equation by its name), the closer names no removed local: t = 0 gives
x^0 (core, and a case).
\begin{stmt}
\sit{cases}{t\in \N,\ x,y\in \Q_{>0},\ t \le 1,\ y \ge 1,\ x \le y:\ x^t \le y}
\end{stmt}
\end{lemma}

\begin{proof}
\begin{plines}
\pline{cases}{Case \((\mathrm{c}_{1})\) \(t = 0\): the goal follows by expanding. Case \((\bar{\mathrm{c}}_{1})\) \(t \ne 0\): \textup{(1)}~\(t = 1\) from the facts at hand. The goal follows by expanding.}
\end{plines}
\end{proof}

\begin{lemma}[{a lower bound of a binomial coefficient}]\label{lem:zchoose}
The word choose bounds a binomial coefficient from below, (m + 1)^d <= binom(m + d, d) d! (Mathlib's
Nat.pow_le_choose), the factorial written as the product 1 * ... * d; also inside with (with), then times a power by
mono (mono, chain: L^d <= q^d d! binom(m + d, d) from L <= q (m + 1)).
\begin{stmt}
\sit{ge}{m,d\in \N:\ (\textstyle\prod_{i=1}^d i) \tbinom{m+d}{d} \ge (m+1)^d}
\sit{with}{m,d\in \N,\ x\in \Q_{\ge0},\ x \le (m+1)^d:\ x \le \tbinom{m+d}{d}\cdot (\textstyle\prod_{i=1}^d i)}
\sit{mono}{m,d\in \N,\ q\in \Q_{>0},\ (m+1)^d \le (\textstyle\prod_{i=1}^d i) \tbinom{m+d}{d}:\ q^d\cdot (m+1)^d \le q^d\cdot ((\textstyle\prod_{i=1}^d i) \tbinom{m+d}{d})}
\sit{chain}{m,d\in \N,\ q,L\in \Q_{>0},\ L \le q\cdot (m+1):\ L^d \le q^d\cdot ((\textstyle\prod_{i=1}^d i) \tbinom{m+d}{d})}
\end{stmt}
\end{lemma}

\begin{proof}
\begin{plines}
\pline{ge}{The goal follows by the bound of the binomial coefficient.}
\pline{with}{\textup{(1)}~\((m+1)^d \le \tbinom{m+d}{d}\cdot (\textstyle\prod_{i=1}^d i)\) by the bound of the binomial coefficient. The goal follows from the facts at hand.}
\pline{mono}{The goal follows by monotonicity.}
\pline{chain}{\textup{(1)}~\((m+1)^d \le (\textstyle\prod_{i=1}^d i) \tbinom{m+d}{d}\) by the bound of the binomial coefficient. \textup{(2)}~\(L^d \le (q\cdot (m+1))^d\) by monotonicity. \textup{(3)}~\(q^d\cdot (m+1)^d \le q^d\cdot ((\textstyle\prod_{i=1}^d i) \tbinom{m+d}{d})\) by monotonicity. The goal follows from the facts at hand.}
\end{plines}
\end{proof}

\begin{lemma}[{the weak compositions in a box}]\label{lem:zsimplex}
The word simplex counts the lattice points of a box [0, c] of NN^r whose coordinates add up to m, every c_(k) at
least m: the weak compositions of m into r = d + 1 parts, binom(m + d, d) (Mathlib's
Finset.card_finsuppAntidiag_nat_eq_choose), either side (flip, with binom(m + d, m)).
\begin{stmt}
\sit{flip}{r,m,d\in \N,\ c\in \mathrm{Grid}(r),\ r = d+1,\ (\forall k:\ c_k \ge m):\ \tbinom{m+d}{m} = \textstyle\sum_{t\le c} \mathbf 1[\textstyle\sum_k t_k = m]}
\end{stmt}
\end{lemma}

\begin{proof}
\begin{plines}
\pline{flip}{The goal follows by the count of the lattice points of a simplex.}
\end{plines}
\end{proof}

\begin{lemma}[{the dyadic exponent of a number in (0, 1]}]\label{lem:zdyadic}
The word dyadic gives the exponent of a number w in (0, 1] on the scale of the powers of a numeral b >= 2,
1 / b^(z + 1) < w <= 1 / b^z (Mathlib's exists_nat_pow_near_of_lt_one at 1 / b), its bounds 0 < w and w <= 1 by V.
\begin{stmt}
\sit{three}{w\in \Q_{\ge0},\ 0 < w,\ w \le 1:\ \exists z\in \N:\ 1/3^{z+1} < w,\ w \le 1/3^z}
\sit{prod}{w,x\in \Q_{>0},\ w \le 1,\ x \le 1/2:\ \exists z\in \N:\ 1/2^{z+1} < w x,\ w x \le 1/2^z}
\sit{coord}{n\in \N,\ w\in \mathrm{Point}(n),\ (\forall i:\ 0 < w_i),\ (\forall i:\ w_i \le 1):\ \forall i:\ \exists z\in \N:\ 1/2^{z+1} < w_i,\ w_i \le 1/2^z}
\end{stmt}
\end{lemma}

\begin{proof}
\begin{plines}
\pline{three}{The goal follows by the dyadic interval of a mass.}
\pline{prod}{The goal follows by the dyadic interval of a mass.}
\pline{coord}{The goal follows by the dyadic interval of a mass.}
\end{plines}
\end{proof}

\begin{lemma}[{a sum of r terms against r times a bound}]\label{lem:zsumconst}
sum[W] reads a side r * t (or t * r), r the length of the tuples that the index of the sum on the other side reads
and t free of that index, as the sum of r terms t, then compares the two sums term by term (le, ge, flip, eq; rows:
the index over the rows of a tuple of tuples), in the type of the relation (nat: a sum of natural numbers); a side
t^r against a product over coordinates, the same way (prod).
\begin{stmt}
\sit{le}{r\in \N,\ x\in \Q_{\ge0},\ F\in \mathrm{Point}(r),\ (\forall j:\ F_j \le x):\ \textstyle\sum_j F_j \le r x}
\sit{ge}{r\in \N,\ x\in \Q_{\ge0},\ F\in \mathrm{Point}(r),\ (\forall j:\ x \le F_j):\ \textstyle\sum_j F_j \ge x r}
\sit{flip}{r\in \N,\ x\in \Q_{\ge0},\ F\in \mathrm{Point}(r),\ (\forall j:\ F_j \le x):\ r x \ge \textstyle\sum_j F_j}
\sit{eq}{r\in \N,\ x\in \Q_{\ge0},\ F\in \mathrm{Point}(r),\ (\forall j:\ F_j = x):\ \textstyle\sum_j F_j = r x}
\sit{rows}{r,n\in \N,\ x,y\in \Q_{>0},\ A\in \mathrm{Rows}(r,n),\ (\forall j:\ y^{\textstyle\sum_i A_{ji}} \le x):\ \textstyle\sum_j y^{\textstyle\sum_i A_{ji}} \le r x}
\sit{nat}{n\in \N,\ b\in \mathrm{Grid}(n),\ (\forall i:\ b_i = 1):\ \textstyle\sum_i b_i = n\cdot 1}
\sit{prod}{n,N\in \N,\ x\in \mathrm{Weight}(n),\ (\forall i:\ x_i = 1/2^{N+1}):\ \textstyle\prod_i x_i = (1/2^{N+1})^n}
\end{stmt}
\end{lemma}

\begin{proof}
\begin{plines}
\pline{le}{The goal follows term by term.}
\pline{ge}{The goal follows term by term.}
\pline{flip}{The goal follows term by term.}
\pline{eq}{The goal follows term by term.}
\pline{rows}{The goal follows term by term.}
\pline{nat}{The goal follows term by term.}
\pline{prod}{The goal follows term by term.}
\end{plines}
\end{proof}

\begin{lemma}[{a loss at a member of a tuple of points}]\label{lem:zalias}
A function on points read at a member of a tuple of tuples: the type of the member is the type of the points once
the type rows are unfolded (Forms(r, n) of tuples of QQ>=0, Point(n) := (QQ>=0)^n).
\begin{stmt}
\sit{member}{r,n\in \N,\ e\in \Q_{>0},\ K\in \mathrm{Loss}(n),\ F\in \mathrm{Forms}(r,n),\ (\forall w\in \mathrm{Point}(n):\ K(w) \le e):\ \forall k:\ K(F_k) \le e}
\end{stmt}
\end{lemma}

\begin{proof}
\begin{plines}
\pline{member}{The goal follows by logic and by logic.}
\end{plines}
\end{proof}

\begin{lemma}[{the facts at the index of a sum}]\label{lem:zsumat}
At the index of a sum compared term by term: a natural term with a numeral read in QQ has its numerals cast at the
leaves, which are dropped (cast, castprod); a fact over the points of the box whose premise is that the point lies
in the box has the premise discharged by the range of the index (boxprem); a fact over every natural number is read
at the index of a sum over a range (range0).
\begin{stmt}
\sit{cast}{k,j,O\in \N,\ A\in \mathrm{Grid}(k),\ F\in \mathrm{Point}(k),\ (\forall \iota:\ F_{\iota} \le A_{\iota}\cdot (j+O)+1):\ \textstyle\sum_{\iota} F_{\iota} \le \textstyle\sum_{\iota} (A_{\iota}\cdot (j+O)+1)}
\sit{castprod}{k,j\in \N,\ A\in \mathrm{Grid}(k),\ F\in \mathrm{Point}(k),\ (\forall \iota:\ F_{\iota} \le A_{\iota} j+1):\ \textstyle\prod_{\iota} F_{\iota} \le \textstyle\prod_{\iota} (A_{\iota} j+1)}
\sit{boxprem}{n\in \N,\ b\in \mathrm{Grid}(n),\ x\in \Q_{\ge0},\ (\forall u\in \mathrm{Grid}(n):\ (\forall i:\ u_i \le b_i)\Rightarrow \mathbf 1[\textstyle\sum_i u_i \le 2] \le x):\ \textstyle\sum_{u\le b} \mathbf 1[\textstyle\sum_i u_i \le 2] \le \textstyle\sum_{u\le b} x}
\sit{range0}{J\in \N,\ x\in \Q_{\ge0},\ c_0,\dots,c_J\in \Q_{\ge0},\ (\forall j\in \N:\ c_j \le x):\ \textstyle\sum_{j=0}^J c_j \le \textstyle\sum_{j=0}^J x}
\end{stmt}
\end{lemma}

\begin{proof}
\begin{plines}
\pline{cast}{The goal follows term by term, from the facts at hand.}
\pline{castprod}{The goal follows term by term, from the facts at hand.}
\pline{boxprem}{The goal follows term by term, from the facts at hand.}
\pline{range0}{The goal follows term by term, from the facts at hand.}
\end{plines}
\end{proof}

\begin{lemma}[{two brackets compared by an implication}]\label{lem:zivle}
An inequality of two Iverson brackets by logic is the implication of their conditions, proved by the word (as the
equation of two brackets is their equivalence): one bracket (core, ge), with a condition that cannot hold (empty),
and term by term in a sum over the rows (sum).
\begin{stmt}
\sit{ge}{p,q\in \Q_{\ge0},\ (p < 1\Rightarrow q < 1):\ \mathbf 1[q < 1] \ge \mathbf 1[p < 1]}
\sit{empty}{p\in \Q_{\ge0},\ p \ge 1:\ \mathbf 1[p < 1] \le \mathbf 1[1 \le 0]}
\sit{sum}{N,n\in \N,\ l\in \mathrm{Forms}(N,n),\ w\in \mathrm{Point}(n),\ e\in \mathrm{Rows}(N,n),\ (\forall k:\ (\forall i:\ l_{ki} < w_i)\Rightarrow \forall i:\ e_{ki} \le 1):\ \textstyle\sum_k \mathbf 1[\forall i:\ l_{ki} < w_i] \le \textstyle\sum_k \mathbf 1[\forall i:\ e_{ki} \le 1]}
\end{stmt}
\end{lemma}

\begin{proof}
\begin{plines}
\pline{ge}{The goal follows by logic.}
\pline{empty}{The goal follows by logic and from the facts at hand.}
\pline{sum}{The goal follows term by term, by logic.}
\end{plines}
\end{proof}

\begin{lemma}[{a tuple of witnesses}]\label{lem:zchoice}
The word choice takes a witness at each coordinate as one tuple (Classical.choose): from forall i : exists x in S :
P(i, x), S one of NN, QQ>=0, QQ>0, the tuple z with forall i : P(i, z_(i)) (nat, rat), over the rows and the
coordinates of a tuple of tuples too (rows); with the dyadic exponents of the coordinates of a point (with).
\begin{stmt}
\sit{nat}{n\in \N,\ w\in \mathrm{Point}(n),\ (\forall i:\ \exists x\in \N:\ 1/2^{x+1} < w_i,\ w_i \le 1/2^x):\ \exists z\in \mathrm{Grid}(n):\ \forall i:\ 1/2^{z_i+1} < w_i,\ w_i \le 1/2^{z_i}}
\sit{rat}{n\in \N,\ w\in \mathrm{Point}(n),\ (\forall i:\ \exists x\in \Q_{\ge0}:\ w_i \le x,\ x \le w_i+1):\ \exists z\in \mathrm{Point}(n):\ \forall i:\ w_i \le z_i,\ z_i \le w_i+1}
\sit{rows}{N,n\in \N,\ l\in \mathrm{Forms}(N,n),\ (\forall k:\ \forall i:\ \exists x\in \N:\ l_{ki} < x):\ \exists z\in \mathrm{Rows}(N,n):\ \forall k:\ \forall i:\ l_{ki} < z_{ki}}
\sit{with}{n\in \N,\ w\in \mathrm{Point}(n),\ (\forall i:\ 0 < w_i),\ (\forall i:\ w_i \le 1):\ \exists z\in \mathrm{Grid}(n):\ \forall i:\ 1/2^{z_i+1} < w_i,\ w_i \le 1/2^{z_i}}
\end{stmt}
\end{lemma}

\begin{proof}
\begin{plines}
\pline{nat}{The goal follows by a witness at each coordinate.}
\pline{rat}{The goal follows by a witness at each coordinate.}
\pline{rows}{The goal follows by a witness at each coordinate.}
\pline{with}{\textup{(1)}~\(\forall i:\ \exists x\in \N:\ 1/2^{x+1} < w_i,\ w_i \le 1/2^x\) by the dyadic interval of a mass. The goal follows by a witness at each coordinate.}
\end{plines}
\end{proof}

\begin{lemma}[{the lattice points of a cut box as a list}]\label{lem:zenum}
The word enum lists the lattice points of a box [0, b] of NN^n (or [0, c]^n) that meet a row R: a tuple of N
tuples e, every e_(k) in the cut box, every point of the cut box some e_(k), and each point at most once (Mathlib's
Finset.equivFin of the filtered box); units lists the unit vectors. The word reindex reads a sum over such a list
as the sum over the box of iverson(R) times the term (Finset.sum_nbij), from the three facts of enum in scope.
\begin{stmt}
\sit{units}{n\in \N:\ \exists N\in \N:\ \exists e\in \mathrm{Rows}(N,n):\ (\forall k:\ (\forall i:\ e_{ki} \le 1),\ \textstyle\sum_i e_{ki} = 1),\ (\forall u\in \mathrm{Grid}(n):\ (\forall i:\ u_i \le 1),\ \textstyle\sum_i u_i = 1\Rightarrow \exists k:\ \forall i:\ e_{ki} = u_i),\ \forall u\in \mathrm{Grid}(n):\ \textstyle\sum_k \mathbf 1[\forall i:\ e_{ki} = u_i] \le 1}
\sit{reindex}{n,N\in \N,\ b\in \mathrm{Grid}(n),\ e\in \mathrm{Rows}(N,n),\ y\in \Q_{>0},\ (\forall k:\ (\forall i:\ e_{ki} \le b_i),\ \textstyle\sum_i e_{ki} \le 2),\ (\forall u\in \mathrm{Grid}(n):\ (\forall i:\ u_i \le b_i),\ \textstyle\sum_i u_i \le 2\Rightarrow \exists k:\ \forall i:\ e_{ki} = u_i),\ (\forall u\in \mathrm{Grid}(n):\ \textstyle\sum_k \mathbf 1[\forall i:\ e_{ki} = u_i] \le 1):\ \textstyle\sum_k y^{\textstyle\sum_i e_{ki}} = \textstyle\sum_{u\le b} \mathbf 1[\textstyle\sum_i u_i \le 2] y^{\textstyle\sum_i u_i}}
\sit{with}{n\in \N,\ b\in \mathrm{Grid}(n),\ y\in \Q_{>0}:\ \exists N\in \N:\ \exists e\in \mathrm{Rows}(N,n):\ \textstyle\sum_k y^{\textstyle\sum_i e_{ki}} = \textstyle\sum_{u\le b} \mathbf 1[\textstyle\sum_i u_i \le 2] y^{\textstyle\sum_i u_i}}
\end{stmt}
\end{lemma}

\begin{proof}
\begin{plines}
\pline{units}{The goal follows by the lattice points of the box.}
\pline{reindex}{The goal follows by the sum over the listed points read as the sum over the box.}
\pline{with}{\textup{(1)}~\(\exists N\in \N:\ \exists e\in \mathrm{Rows}(N,n):\ (\forall k:\ (\forall i:\ e_{ki} \le b_i),\ \textstyle\sum_i e_{ki} \le 2),\ (\forall u\in \mathrm{Grid}(n):\ (\forall i:\ u_i \le b_i),\ \textstyle\sum_i u_i \le 2\Rightarrow \exists k:\ \forall i:\ e_{ki} = u_i),\ \forall u\in \mathrm{Grid}(n):\ \textstyle\sum_k \mathbf 1[\forall i:\ e_{ki} = u_i] \le 1\) by the lattice points of the box. Take \(N=N, e=e\). The goal follows by the sum over the listed points read as the sum over the box.}
\end{plines}
\end{proof}

\begin{fact}[{the floor, Mathlib \texttt{Nat.\allowbreak floor\_\allowbreak le}, \texttt{Nat.\allowbreak lt\_\allowbreak floor\_\allowbreak add\_\allowbreak one}}]\label{fact:zfloor}
A mathlib: row may name several lemmas that together prove the cited statement, here the floor. logic and with
destruct an existential over a number set (QQ>0, QQ>=0, NN, NN>=k), the sign of its witness a fact, and logic
destructs the existentials among the hypotheses in scope, the rows of a relation too.
\begin{stmt}
\sit{core}{x\in \Q_{\ge0}:\ \exists k\in \N:\ k \le x,\ x < k+1}
\end{stmt}
\end{fact}

\begin{stmt}
\sit{relation}{c\ \text{zpos}\ L:\iff c\in \Q_{\ge0},\ L\in \N,\ \exists d\in \Q_{>0}:\ \exists M\in \N_{\ge1}:\ d \le c,\ L \le M}
\end{stmt}

\begin{lemma}[{the floor applied, number witnesses destructed}]\label{lem:zfloors}
\begin{stmt}
\sit{apply}{x\in \Q_{\ge0}:\ \exists k\in \N:\ k \le x,\ x < k+1}
\sit{relation}{c\in \Q_{\ge0},\ L\in \N,\ c\ \text{zpos}\ L:\ \exists e\in \Q_{>0}:\ \exists N\in \N_{\ge1}:\ e \le c+1,\ L \le N+1}
\sit{given}{c\in \Q_{\ge0},\ (\exists d\in \Q_{>0}:\ d \le c):\ \exists e\in \Q_{>0}:\ e \le 2 c}
\end{stmt}
\end{lemma}

\begin{proof}
\begin{plines}
\pline{apply}{The goal follows by \cref{fact:zfloor}(core).}
\pline{relation}{By logic: take \(e=d, N=M\). Each part of the goal follows from the facts at hand.}
\pline{given}{By logic: take \(e=2 d\). The goal follows from the facts at hand.}
\end{plines}
\end{proof}

\begin{lemma}[{a region named in a proof}]\label{lem:zregion}
take[R := {w in X : P}; W] names a region by a predicate, in a proof only (a statement takes a region variable with
its characterizing hypothesis). R, of a type row set(X), comes with two facts that
logic applies, forall w in R : P and forall w in X : P => exists v in R : forall i : v_(i) = w_(i) (the membership
of a point as the rows write it); a reference takes R as data. Here a box in a larger one, by a rule on regions.
\begin{stmt}
\sit{cover}{n\in \N,\ B,B'\in \mathrm{Region}(n),\ (\forall w\in B:\ \exists v\in B':\ \forall i:\ v_i = w_i),\ (\forall w\in B':\ \forall i:\ w_i \le 2):\ \forall w\in B:\ \forall i:\ w_i \le 2}
\sit{box}{n\in \N:\ \forall w\in \mathrm{Point}(n):\ (\forall i:\ w_i \le 1)\Rightarrow \forall i:\ w_i \le 2}
\end{stmt}
\end{lemma}

\begin{proof}
\begin{plines}
\pline{cover}{The goal follows by logic and from the facts at hand.}
\pline{box}{Take \(R=\{w \in \mathrm{Point}(n):\ \forall i:\ w_i \le 1\}, R'=\{w \in \mathrm{Point}(n):\ \forall i:\ w_i \le 2\}\). \textup{(1)}~\(\forall u\in R:\ \exists v\in R':\ \forall i:\ v_i = u_i\). By logic: \textup{(2)}~\(\forall i:\ u_i \le 2\) from the facts at hand. The goal follows by logic. This proves \textup{(1)}. \textup{(3)}~\(\forall u\in R:\ \forall i:\ u_i \le 2\) by \cref{lem:zregion}(cover). The goal follows by logic and from the facts at hand.}
\end{plines}
\end{proof}

\begin{lemma}[{a case split}]\label{lem:zsplit}
cases[P; W1; W2] splits on the clause P: W1 proves the goal with P as a fact, W2 with its negation as a fact,
pushed in (a relation flipped, forall and exists swapped, a conjunction under not an implication; not n = 0 for a
natural number n is 1 <= n). Here a coordinate zero or at least 1, a point with a large coordinate or none, and two
orders.
\begin{stmt}
\sit{zero}{n\in \N,\ a\in \mathrm{Grid}(n),\ y\in \mathrm{Point}(n),\ x\in \Q_{\ge0},\ (\forall i:\ 1 \le a_i\Rightarrow y_i \le x),\ (\forall i:\ a_i = 0\Rightarrow y_i \le x):\ \forall i:\ y_i \le x}
\sit{shell}{n\in \N,\ a\in \mathrm{Grid}(n),\ s,y\in \mathrm{Point}(n),\ K\in \Q_{\ge0},\ (\forall i:\ 1 \le a_i\Rightarrow s_i \le 2^{a_i} y_i\Rightarrow K \le 1),\ ((\forall i:\ 1 \le a_i\Rightarrow 2^{a_i} y_i < s_i)\Rightarrow K \le 1):\ K \le 1}
\sit{order}{x,y\in \Q_{\ge0},\ (x \le y\Rightarrow x \le 2),\ (y < x\Rightarrow x \le 2+y),\ y \le 1:\ x \le 3}
\end{stmt}
\end{lemma}

\begin{proof}
\begin{plines}
\pline{zero}{Case \((\mathrm{c}_{1})\) \(a_i = 0\): the goal follows by logic and from the facts at hand. Case \((\bar{\mathrm{c}}_{1})\) \(a_i \ne 0\): the goal follows by logic and from the facts at hand.}
\pline{shell}{Case \((\mathrm{c}_{1})\) \(\exists i:\ 1 \le a_i,\ s_i \le 2^{a_i} y_i\): the goal follows by logic and from the facts at hand. Case \((\bar{\mathrm{c}}_{1})\) \(\lnot(\exists i:\ 1 \le a_i,\ s_i \le 2^{a_i} y_i)\): the goal follows by logic and from the facts at hand.}
\pline{order}{Case \((\mathrm{c}_{1})\) \(x \le y\): the goal follows by logic and from the facts at hand. Case \((\bar{\mathrm{c}}_{1})\) \(x > y\): the goal follows by logic and from the facts at hand.}
\end{plines}
\end{proof}

\begin{lemma}[{a point of a region by its characterization}]\label{lem:zmember}
A point of a region as the rows write it, exists v in B : forall i : v_(i) = w_(i), is the point w itself (v is
identified with w), so a rule over B applies at w. A rule that the proof brings in (a fact of a region named by
take, a step of with) is applied to the hypotheses already in scope too.
\begin{stmt}
\sit{bound}{n\in \N,\ B\in \mathrm{Region}(n),\ K\in \mathrm{Loss}(n),\ (\forall w\in B:\ K(w) \le 1),\ u\in \mathrm{Point}(n),\ (\exists v\in B:\ \forall i:\ v_i = u_i):\ K(u) \le 1}
\sit{named}{n\in \N,\ K\in \mathrm{Loss}(n),\ (\forall w\in \mathrm{Point}(n):\ (\forall i:\ w_i \le 1)\Rightarrow K(w) \le 1),\ u\in \mathrm{Point}(n),\ (\forall i:\ u_i \le 1):\ 2 K(u) \le 2}
\end{stmt}
\end{lemma}

\begin{proof}
\begin{plines}
\pline{bound}{The goal follows by logic and from the facts at hand.}
\pline{named}{Take \(R=\{w \in \mathrm{Point}(n):\ \forall i:\ w_i \le 1\}\). \textup{(1)}~\(\forall w\in R:\ K(w) \le 1\) by logic and from the facts at hand. The goal follows by logic and from the facts at hand.}
\end{plines}
\end{proof}

\begin{lemma}[{rules over numbers}]\label{lem:zrules}
logic reads a binder over a number set (its bound, NN>=1 or QQ>0, by arithmetic) and a premise of several
conjuncts (each matched with a fact alone, passed together); a universal goal over several names is introduced
name by name.
\begin{stmt}
\sit{rule}{\lambda\in \Q_{>0},\ L_0\in \N,\ (\forall L,\kappa\in \N:\ L \ge L_0,\ \kappa \le \lambda L,\ \lambda L < \kappa+1\Rightarrow \exists N\in \N_{\ge1}:\ \kappa \le N):\ \forall L,\kappa\in \N:\ L \ge L_0+1,\ \kappa \le \lambda L,\ \lambda L < \kappa+1\Rightarrow \exists N\in \N_{\ge1}:\ \kappa \le N}
\sit{guard}{(\forall L\in \N_{\ge1}:\ L \le 2\Rightarrow L L \le 4):\ \forall L\in \N:\ L = 1\Rightarrow L L \le 4}
\end{stmt}
\end{lemma}

\begin{proof}
\begin{plines}
\pline{rule}{\textup{(1)}~\(L \ge L_0\) from the facts at hand. By logic: take \(N=N\). The goal follows from the facts at hand.}
\pline{guard}{\textup{(1)}~\(L \le 2\) from the facts at hand. The goal follows by logic and from the facts at hand.}
\end{plines}
\end{proof}

\begin{lemma}[{a key of two words}]\label{lem:zkeys}
A reference names a row whose key has several words: lem:y(two words), the Lean name with underscores.
\begin{stmt}
\sit{first bound}{x\in \Q_{\ge0},\ x \le 1:\ x \le 2}
\sit{use}{x\in \Q_{\ge0},\ x \le 1:\ x \le 3}
\end{stmt}
\end{lemma}

\begin{proof}
\begin{plines}
\pline{first bound}{The goal follows from the facts at hand.}
\pline{use}{The goal follows by \cref{lem:zkeys}(first bound).}
\end{plines}
\end{proof}

\begin{lemma}[{the closer on powers, quotients, maxima and natural goals}]\label{lem:zclose}
The closer, last: the conjunctions among the facts split, an exponent with a sum or a factor 0 or 1 distributed
((l + 1) * a = l * a + a, then the power split), a quotient by a positive term cleared (k <= x / D gives
k * D <= x), a maximum split by cases, a natural goal read in QQ (k <= j as k < j + 1), the relations between
natural numbers of every fact read in QQ, a sign 0 <= t by positivity, and for an equation the equations of the
facts substituted before the quotients are cleared.
\begin{stmt}
\sit{power}{n,l\in \N,\ a\in \mathrm{Grid}(n),\ y,u,\mathrm{lb}\in \mathrm{Point}(n),\ \mathrm{sb}\in \mathrm{Weight}(n),\ (\forall i:\ \mathrm{lb}_i \le y_i,\ 2^{(l+1) a_i} y_i \le \mathrm{lb}_i+\mathrm{sb}_i),\ (\forall i:\ u_i = 2^{a_i} y_i):\ \forall i:\ \mathrm{lb}_i \le u_i,\ 2^{l a_i} u_i \le \mathrm{lb}_i+\mathrm{sb}_i}
\sit{zero}{n\in \N,\ a\in \mathrm{Grid}(n),\ y,\mathrm{lb}\in \mathrm{Point}(n),\ \mathrm{sb}\in \mathrm{Weight}(n),\ (\forall i:\ \mathrm{lb}_i \le y_i,\ 2^{0 a_i} y_i \le \mathrm{lb}_i+\mathrm{sb}_i):\ \forall i:\ y_i \le \mathrm{lb}_i+\mathrm{sb}_i}
\sit{level}{D\in \N_{\ge1},\ L,K\in \N:\ \exists l_1\in \N:\ L+K \le l_1 D,\ l_1 D \le L+K+D}
\sit{deficit}{D\in \N_{\ge1},\ L,K\in \N,\ K+1 \le L:\ \exists g\in \N:\ g D+K+1 \le L,\ L \le g D+K+D}
\sit{scale}{n\in \N,\ a\in \mathrm{Grid}(n),\ w,w'\in \mathrm{Point}(n),\ h\in \mathrm{Weight}(n),\ (\forall i:\ w'_i = h_i w_i),\ (\forall i:\ h_i = 1/2^{a_i}):\ \forall i:\ 2^{a_i} w'_i = w_i}
\sit{cast}{c\in \Q_{>0},\ J,L\in \N,\ J > 4/c,\ L \ge 2 J:\ 4/c < L+1}
\end{stmt}
\end{lemma}

\begin{proof}
\begin{plines}
\pline{power}{Each part of the goal follows from the facts at hand.}
\pline{zero}{The goal follows from the facts at hand.}
\pline{level}{\textup{(1)}~\(\exists k\in \N:\ k \le (L+K)/D,\ (L+K)/D < k+1\) by \cref{fact:zfloor}(core) (\((L+K)/D \ge 0\) from the facts at hand). Take \(l_1=k+1\). Each part of the goal follows from the facts at hand.}
\pline{deficit}{\textup{(1)}~\(\exists k\in \N:\ k \le \max(L-(K+1),0)/D,\ \max(L-(K+1),0)/D < k+1\) by \cref{fact:zfloor}(core) (\(\max(L-(K+1),0)/D \ge 0\) from the facts at hand). Take \(g=k\). Each part of the goal follows from the facts at hand.}
\pline{scale}{The goal follows from the facts at hand.}
\pline{cast}{\textup{(1)}~\(L \ge J\) from the facts at hand. The goal follows from the facts at hand.}
\end{plines}
\end{proof}

\begin{lemma}[{a region by its characterization}]\label{lem:zchar}
logic reads an equivalence that ends a hypothesis as its two implications, and takes a membership of a point for
the membership as the rows write it (exists v in C : forall i : v_(i) = w_(i)): a region given by its
characterization, read both ways.
\begin{stmt}
\sit{forward}{n\in \N,\ C\in \mathrm{Region}(n),\ K\in \mathrm{Loss}(n),\ (\forall y\in \mathrm{Point}(n):\ (\exists v\in C:\ \forall i:\ v_i = y_i)\iff (\forall i:\ y_i \le 1)),\ (\forall y\in \mathrm{Point}(n):\ (\forall i:\ y_i \le 1)\Rightarrow K(y) \le 1):\ \forall w\in C:\ K(w) \le 1}
\sit{back}{n\in \N,\ C\in \mathrm{Region}(n),\ K\in \mathrm{Loss}(n),\ (\forall y\in \mathrm{Point}(n):\ (\exists v\in C:\ \forall i:\ v_i = y_i)\iff (\forall i:\ y_i \le 1)),\ (\forall w\in C:\ K(w) \le 1),\ u\in \mathrm{Point}(n),\ (\forall i:\ u_i \le 1):\ K(u) \le 1}
\end{stmt}
\end{lemma}

\begin{proof}
\begin{plines}
\pline{forward}{The goal follows by logic and from the facts at hand.}
\pline{back}{\textup{(1)}~\(\forall i:\ u_i \le 1\) from the facts at hand. The goal follows by logic and from the facts at hand.}
\end{plines}
\end{proof}

\begin{lemma}[{one word on a relation goal}]\label{lem:zone}
One word on a relation goal of several rows: a reference proves it whole; V or logic takes the fact in scope that
states it (a step of with), else closes its fields from the facts.
\begin{stmt}
\sit{site}{q\in \Q_{>0},\ t\in \Q_{\ge0},\ q \le t,\ t \le q+q:\ q\ \text{zsite}\ t}
\sit{again}{q\in \Q_{>0},\ t\in \Q_{\ge0},\ q \le t,\ t \le q+q:\ q\ \text{zsite}\ t}
\sit{step}{q\in \Q_{>0},\ t\in \Q_{\ge0},\ q \le t,\ t \le q+q:\ q\ \text{zsite}\ t}
\end{stmt}
\end{lemma}

\begin{proof}
\begin{plines}
\pline{site}{Each part of the goal follows from the facts at hand.}
\pline{again}{The goal follows by \cref{lem:zone}(site).}
\pline{step}{\textup{(1)}~\(q\ \text{zsite}\ t\) by \cref{lem:zone}(site). The goal follows from the facts at hand.}
\end{plines}
\end{proof}

logic, with its default word: a membership of a union is split by cases, each case a logic of its own with the
membership in its region (the rules over that region apply at the point), closed by the goal as a fact, by
arithmetic (exfalso; linarith) or by tauto; a membership known already is not drawn again, so that a cycle of
inclusions ends. A relation fact (a step of with, a witness) gives its fields that are rules.

\begin{stmt}
\sit{relation}{K\ \text{zcapped}\ L\ \text{zon}\ B:\iff n,L\in \N,\ K\in \mathrm{Loss}(n),\ B\in \mathrm{Region}(n),\ \forall w\in B:\ 2^L K(w) \le 1\Rightarrow \exists k\in \N:\ K(w) \le k}
\end{stmt}

\begin{lemma}[{a membership of a union, split}]\label{lem:zunion}
\begin{stmt}
\sit{union}{n,L\in \N,\ K\in \mathrm{Loss}(n),\ B,B_1,B_2\in \mathrm{Region}(n),\ B \subseteq B_1\cup B_2,\ K\ \text{zcapped}\ L\ \text{zon}\ B_1,\ (\forall w\in B_2:\ 1 < 2^L K(w)):\ K\ \text{zcapped}\ L\ \text{zon}\ B}
\sit{cycle}{n,L\in \N,\ K\in \mathrm{Loss}(n),\ B,B_1,B_2\in \mathrm{Region}(n),\ B \subseteq B_1\cup B_2,\ B_2 \subseteq B,\ K\ \text{zcapped}\ L\ \text{zon}\ B_1,\ (\forall w\in B_2:\ 1 < 2^L K(w)):\ K\ \text{zcapped}\ L\ \text{zon}\ B}
\end{stmt}
\end{lemma}

\begin{proof}
\begin{plines}
\pline{union}{\textup{(1)}~\(K\ \text{zcapped}\ L\ \text{zon}\ B_1\) by logic. The goal follows by logic.}
\pline{cycle}{\textup{(1)}~\(K\ \text{zcapped}\ L\ \text{zon}\ B_1\) by logic. The goal follows by logic.}
\end{plines}
\end{proof}

\begin{lemma}[{the parts of a premise}]\label{lem:zparts}
A premise of several parts stays a fact part by part when logic destructs it (an existential among them too), so
a rule whose premise has the same parts applies; a universal over the coordinates inside a conjunction is
instantiated at the index the goal introduces; the closer substitutes an equation between a natural value and a
numeral (a_(i) = 0) before it reads the powers. A relation that a fact states as it is closes by that fact.
\begin{stmt}
\sit{low}{n\in \N,\ a\in \mathrm{Grid}(n),\ \mathrm{lb},y\in \mathrm{Point}(n),\ \mathrm{sb}\in \mathrm{Weight}(n),\ S\in \mathrm{Region}(n),\ c_0\in \Q_{>0},\ K\in \mathrm{Loss}(n),\ (\forall v\in \mathrm{Point}(n):\ (\exists u\in S:\ \forall i:\ u_i = v_i)\iff ((\forall i:\ \mathrm{lb}_i \le v_i \le \mathrm{lb}_i+\mathrm{sb}_i),\ (\exists i:\ 1 \le a_i,\ \mathrm{sb}_i \le 2^{a_i} v_i))),\ (\forall w\in S:\ c_0 \le K(w)):\ (\forall i:\ \mathrm{lb}_i \le y_i \le \mathrm{lb}_i+\mathrm{sb}_i),\ (\exists i:\ 1 \le a_i,\ \mathrm{sb}_i \le 2^{a_i} y_i)\Rightarrow c_0 \le K(y)}
\sit{index}{n\in \N,\ \mathrm{lb},y\in \mathrm{Point}(n),\ \mathrm{sb}\in \mathrm{Weight}(n):\ (\forall i:\ \mathrm{lb}_i \le y_i),\ (\exists i:\ \mathrm{sb}_i \le y_i)\Rightarrow \forall i:\ \mathrm{lb}_i \le y_i+\mathrm{sb}_i}
\sit{zero}{n\in \N,\ a\in \mathrm{Grid}(n),\ u,y,\mathrm{lb}\in \mathrm{Point}(n),\ (\forall i:\ u_i = 2^{a_i} y_i),\ (\forall i:\ y_i \le \mathrm{lb}_i):\ \forall i:\ a_i = 0\Rightarrow u_i \le \mathrm{lb}_i}
\sit{same}{q\in \Q_{>0},\ t\in \Q_{\ge0},\ q \le t,\ t \le q+q:\ \exists r\in \Q_{>0}:\ r\ \text{zsite}\ t}
\end{stmt}
\end{lemma}

\begin{proof}
\begin{plines}
\pline{low}{The goal follows by logic and from the facts at hand.}
\pline{index}{The goal follows from the facts at hand.}
\pline{zero}{The goal follows from the facts at hand.}
\pline{same}{\textup{(1)}~\(q\ \text{zsite}\ t\) from the facts at hand (each part). Take \(r=q\). The goal follows from the facts at hand.}
\end{plines}
\end{proof}

\begin{lemma}[{an existential premise from its parts}]\label{lem:zpack}
An existential premise of a rule (exists i : P_1, P_2) is given by its parts at an index in scope (a destructed
existential, a case of cases); the premises and the facts are matched after the arithmetic of numerals that
push_cast does on the goal ((0 + 1) * a as a). In an equation goal the closer substitutes an equation over every
index (forall i : u_(i) = F) at the index of the goal.
\begin{stmt}
\sit{part}{n\in \N,\ a\in \mathrm{Grid}(n),\ y\in \mathrm{Point}(n),\ \mathrm{sb}\in \mathrm{Weight}(n),\ c_0\in \Q_{>0},\ K\in \mathrm{Loss}(n),\ (\forall v\in \mathrm{Point}(n):\ (\exists i:\ 1 \le a_i,\ \mathrm{sb}_i \le 2^{a_i} v_i)\Rightarrow c_0 \le K(v)):\ (\exists i:\ 1 \le a_i,\ \mathrm{sb}_i \le 2^{(0+1) a_i} y_i)\Rightarrow c_0 \le K(y)}
\sit{case}{n\in \N,\ a\in \mathrm{Grid}(n),\ y\in \mathrm{Point}(n),\ \mathrm{sb}\in \mathrm{Weight}(n),\ c_0\in \Q_{>0},\ K\in \mathrm{Loss}(n),\ (\forall v\in \mathrm{Point}(n):\ (\forall i:\ v_i \le \mathrm{sb}_i),\ (\exists i:\ 1 \le a_i,\ \mathrm{sb}_i \le 2^{a_i} v_i)\Rightarrow c_0 \le K(v)),\ (\forall v\in \mathrm{Point}(n):\ (\forall i:\ 1 \le a_i\Rightarrow 2^{a_i} v_i < \mathrm{sb}_i)\Rightarrow c_0 \le K(v)):\ (\forall i:\ y_i \le \mathrm{sb}_i)\Rightarrow c_0 \le K(y)}
\sit{dilate}{n,l\in \N,\ a\in \mathrm{Grid}(n),\ u,y\in \mathrm{Point}(n),\ \mathrm{sb}\in \mathrm{Weight}(n),\ (\forall i:\ u_i = 2^{a_i} y_i),\ (\exists i:\ \mathrm{sb}_i \le 2^{l a_i} u_i):\ \exists i:\ \mathrm{sb}_i \le 2^{(l+1) a_i} y_i}
\end{stmt}
\end{lemma}

\begin{proof}
\begin{plines}
\pline{part}{The goal follows by logic and from the facts at hand.}
\pline{case}{Case \((\mathrm{c}_{1})\) \(\exists i:\ 1 \le a_i,\ \mathrm{sb}_i \le 2^{a_i} y_i\): the goal follows by logic and from the facts at hand. Case \((\bar{\mathrm{c}}_{1})\) \(\lnot(\exists i:\ 1 \le a_i,\ \mathrm{sb}_i \le 2^{a_i} y_i)\): the goal follows by logic and from the facts at hand.}
\pline{dilate}{By logic: take \(i=i\). \textup{(1)}~\(2^{(l+1) a_i} y_i = 2^{l a_i} u_i\) from the facts at hand. The goal follows from the facts at hand.}
\end{plines}
\end{proof}

A list of one box (k = 1) given by take with a formula that does not read its row: the index of the rows keeps
its range, a universal over the rows is introduced, an existential over them has the only row as its witness,
and a sum or a product over the rows is its term (Fin.sum_univ_one).

\begin{stmt}
\sit{relation}{l\ \text{zbox}\ s\ \text{zin}\ B:\iff k,n\in \N,\ l\in \mathrm{Forms}(k,n),\ s\in \mathrm{Sides}(k,n),\ B\in \mathrm{Region}(n),\ (\forall \iota:\ \forall w\in \mathrm{Point}(n):\ (\forall i:\ l_{\iota,i} \le w_i \le l_{\iota,i}+s_{\iota,i})\Rightarrow \exists v\in B:\ \forall i:\ v_i = w_i),\ \exists \iota:\ \forall i:\ l_{\iota,i} < l_{\iota,i}+s_{\iota,i}}
\end{stmt}

\begin{lemma}[{a list of one box}]\label{lem:zonebox}
\begin{stmt}
\sit{one}{n\in \N,\ \mathrm{lb}\in \mathrm{Point}(n),\ \mathrm{sb}\in \mathrm{Weight}(n),\ B\in \mathrm{Region}(n),\ (\forall w\in \mathrm{Point}(n):\ (\forall i:\ \mathrm{lb}_i \le w_i \le \mathrm{lb}_i+\mathrm{sb}_i)\Rightarrow \exists v\in B:\ \forall i:\ v_i = w_i):\ \exists l\in \mathrm{Forms}(1,n):\ \exists s\in \mathrm{Sides}(1,n):\ l\ \text{zbox}\ s\ \text{zin}\ B,\ \textstyle\sum_{\iota} (\textstyle\prod_i s_{\iota,i}) = \textstyle\prod_i \mathrm{sb}_i}
\end{stmt}
\end{lemma}

\begin{proof}
\begin{plines}
\pline{one}{Take \(l_{\iota,i}=\mathrm{lb}_i, s_{\iota,i}=\mathrm{sb}_i\). The parts of the goal follow: (i) with its parts (i) by logic; (ii) from the facts at hand; (ii) term by term, from the facts at hand.}
\end{plines}
\end{proof}

\begin{lemma}[{quotients and powers in the closer}]\label{lem:zquot}
Quotients and powers in the closer: a quotient by a power of a numeral base at least 1 is at most its numerator,
and times the power it is its numerator; a power whose exponent has a factor 0 is 1; a power against a power of
the same base is monotone in the exponent (tried first); an Iverson bracket is split by cases; V on a chain is V
on each link.
\begin{stmt}
\sit{quot}{n\in \N,\ a\in \mathrm{Grid}(n),\ w,\mathrm{lb}\in \mathrm{Point}(n),\ \mathrm{sb}\in \mathrm{Weight}(n),\ (\forall i:\ w_i \le \mathrm{lb}_i+\mathrm{sb}_i/2^{a_i}):\ \forall i:\ w_i \le \mathrm{lb}_i+\mathrm{sb}_i}
\sit{zero}{n,l\in \N,\ a\in \mathrm{Grid}(n),\ w,\mathrm{lb}\in \mathrm{Point}(n),\ \mathrm{sb}\in \mathrm{Weight}(n),\ (\forall i:\ w_i \le \mathrm{lb}_i+\mathrm{sb}_i/2^{l a_i}):\ \forall i:\ a_i = 0\Rightarrow 2^{l a_i} w_i \le \mathrm{lb}_i+\mathrm{sb}_i}
\sit{cancel}{n,l\in \N,\ a\in \mathrm{Grid}(n),\ w,\mathrm{lb}\in \mathrm{Point}(n),\ \mathrm{sb}\in \mathrm{Weight}(n),\ (\forall i:\ \mathrm{lb}_i = 0),\ (\forall i:\ w_i \le \mathrm{lb}_i+\mathrm{sb}_i/2^{l a_i}):\ \forall i:\ 2^{l a_i} w_i \le \mathrm{lb}_i+\mathrm{sb}_i}
\sit{mono}{k,m,l\in \N,\ k+1 \le m:\ 2^{l\cdot (k+1)} \le 2^{l m}}
\sit{bracket}{x\in \Q_{\ge0}:\ \mathbf 1[x \le 1] \le 1}
\sit{chain}{x,y,z\in \Q_{\ge0},\ x \le y,\ y \le z:\ x \le y \le z}
\end{stmt}
\end{lemma}

\begin{proof}
\begin{plines}
\pline{quot}{The goal follows from the facts at hand.}
\pline{zero}{The goal follows from the facts at hand.}
\pline{cancel}{The goal follows from the facts at hand.}
\pline{mono}{The goal follows from the facts at hand.}
\pline{bracket}{The goal follows from the facts at hand.}
\pline{chain}{The goal follows from the facts at hand.}
\end{plines}
\end{proof}

\begin{lemma}[{an equation of naturals read right to left}]\label{lem:zprod}
An equation of naturals read in QQ from right to left: the power of the sum replaced by the product of the
powers, after the equation e = sum i, a_(i) is substituted.
\begin{stmt}
\sit{core}{n,l\in \N,\ a\in \mathrm{Grid}(n):\ \textstyle\prod_i 2^{l a_i} = 2^{l\cdot (\textstyle\sum_i a_i)}}
\sit{use}{n,l,e\in \N,\ a\in \mathrm{Grid}(n),\ \mathrm{sb}\in \mathrm{Weight}(n),\ e = \textstyle\sum_i a_i:\ (\textstyle\prod_i \mathrm{sb}_i/2^{l a_i})\cdot 2^{l e} = \textstyle\prod_i \mathrm{sb}_i}
\end{stmt}
\end{lemma}

\begin{proof}
\begin{plines}
\pline{use}{\textup{(1)}~\((\textstyle\prod_i \mathrm{sb}_i/2^{l a_i})\cdot (\textstyle\prod_i 2^{l a_i}) = \textstyle\prod_i \mathrm{sb}_i\) term by term, from the facts at hand. \textup{(2)}~\(\textstyle\prod_i 2^{l a_i} = 2^{l\cdot (\textstyle\sum_i a_i)}\) by \cref{lem:zprod}(core). The goal follows from the facts at hand.}
\end{plines}
\end{proof}

\begin{lemma}[{a reference that concludes a universal}]\label{lem:zall}
A reference that concludes a universal over points proves the goal at the point the goal introduced.
\begin{stmt}
\sit{first}{n\in \N,\ K\in \mathrm{Loss}(n),\ c\in \Q_{>0},\ (\forall y\in \mathrm{Point}(n):\ c \le K(y)):\ \forall y\in \mathrm{Point}(n):\ c \le K(y)+c}
\sit{second}{n\in \N,\ K\in \mathrm{Loss}(n),\ c\in \Q_{>0},\ (\forall y\in \mathrm{Point}(n):\ c \le K(y)):\ \forall y\in \mathrm{Point}(n):\ c \le K(y)+c}
\end{stmt}
\end{lemma}

\begin{proof}
\begin{plines}
\pline{first}{The goal follows by logic and from the facts at hand.}
\pline{second}{The goal follows by \cref{lem:zall}(first).}
\end{plines}
\end{proof}

\begin{lemma}[{a universal case clause at the index in scope}]\label{lem:zcase}
A universal case clause holds at the index in scope; in the other case the witness of its negation, an index
named apart (i_2), carries the facts over every index.
\begin{stmt}
\sit{one}{n\in \N,\ w,\mathrm{lb}\in \mathrm{Point}(n),\ \mathrm{sb}\in \mathrm{Weight}(n),\ (\forall i:\ w_i \le \mathrm{lb}_i+\mathrm{sb}_i):\ \forall i:\ w_i \le \mathrm{lb}_i+2 \mathrm{sb}_i}
\end{stmt}
\end{lemma}

\begin{proof}
\begin{plines}
\pline{one}{By logic: Case \((\mathrm{c}_{1})\) \(\forall i:\ w_i \le \mathrm{lb}_i+\mathrm{sb}_i\): the goal follows from the facts at hand. Case \((\bar{\mathrm{c}}_{1})\) \(\lnot(\forall i:\ w_i \le \mathrm{lb}_i+\mathrm{sb}_i)\): the goal follows from the facts at hand.}
\end{plines}
\end{proof}

A name bound in a relation that is free in an argument (C below) is renamed in the instance; take gives it by
its own name.

\begin{stmt}
\sit{relation}{x\ \text{zroom}\ B:\iff x,B\in \Q_{>0},\ x \le B,\ \exists C\in \Q_{>0}:\ C x \le B}
\end{stmt}

\begin{lemma}[{a bound name of a relation free in an argument}]\label{lem:zcap}
\begin{stmt}
\sit{one}{c,C\in \Q_{>0},\ c \le C:\ c\ \text{zroom}\ C}
\end{stmt}
\end{lemma}

\begin{proof}
\begin{plines}
\pline{one}{Part (i): the goal follows from the facts at hand. Part (ii): take \(C=1\). The goal follows from the facts at hand.}
\end{plines}
\end{proof}

\begin{lemma}[{the sign of a product of coordinates in QQ>0}]\label{lem:zprodpos}
The sign of a product of coordinates in QQ>0, alone or over a power, in a side goal of take.
\begin{stmt}
\sit{one}{n,k\in \N,\ \mathrm{sb}\in \mathrm{Weight}(n),\ P\in \Q_{>0},\ P = \textstyle\prod_i \mathrm{sb}_i:\ \exists c\in \Q_{>0}:\ c\cdot 2^k = P}
\end{stmt}
\end{lemma}

\begin{proof}
\begin{plines}
\pline{one}{Take \(c=(\textstyle\prod_i \mathrm{sb}_i)/2^k\). The goal follows from the facts at hand.}
\end{plines}
\end{proof}

\begin{lemma}[{a price made strict by a hypothesis}]\label{lem:progstrict}
A price of nonnegative sign that a strict hypothesis makes strict (0 < kq): the price kq / (1 + eta) is used
from above.
\begin{stmt}
\sit{given}{\mathrm{kq},\eta\in \Q_{\ge0}}
\sit{(1)}{0 < \mathrm{kq}:\ \operatorname{prog}(x\in \Q_{\ge0},\ (1+\eta) x \ge 1:\ \mathrm{kq}\cdot x) = (\mathrm{kq}/(1+\eta),1)}
\end{stmt}
\end{lemma}

\begin{proof}
\begin{plines}
\pline{(1)}{The goal follows by the prices \(\mathrm{kq}/(1+\eta)\) of the rows at the points \(x=1/(1+\eta)\).}
\end{plines}
\end{proof}

\begin{lemma}[{a program of two scales}]\label{lem:zprogrow}
A program row of two scales.
\begin{stmt}
\sit{given}{N\in \Q_{>0}}
\sit{(1)}{\operatorname{prog}(x,y\in \Q_{\ge0},\ 4 x+2 y \ge 1:\ N x+N y) = (N/4,1)}
\end{stmt}
\end{lemma}

\begin{proof}
\begin{plines}
\pline{(1)}{The goal follows by the prices \(N/4\) of the rows at the points \(x=1/4\).}
\end{plines}
\end{proof}

\begin{lemma}[{the reading of a program at a point}]\label{lem:zprogread}
The reading of a program row at a point (at x := a, y := b): its least value is at most the cost there, the point
feasible by the facts; the goal is the reading (1), or follows from it by linear arithmetic (2), also for a row with
a guard (3).
\begin{stmt}
\sit{given}{N,a,b\in \Q_{>0},\ 1 \le 4 a+2 b}
\sit{(1)}{N/4 \le N a+N b}
\sit{(2)}{N \le 4 N a+4 N b}
\end{stmt}
\end{lemma}

\begin{proof}
\begin{plines}
\pline{(1)}{The goal follows by \cref{lem:zprogrow}(1) (at \(x=a, y=b\)).}
\pline{(2)}{The goal follows by \cref{lem:zprogrow}(1) (at \(x=a, y=b\)).}
\end{plines}
\end{proof}

A relation at the level L read from one at the level L' + D: an equation of natural numbers L = t, L a name of
the goal and t free of L, is substituted when logic destructs it (a premise or a witness, both orientations) or
when with states it; the goal then reads t, and the relation is the fact as it is. A fact that binds L again (a
rule over every level, the hypothesis of an induction) keeps its binder, and its number binders are matched to
terms (L := L' + D, rule); the casts of L' + D to QQ are pushed, as a with-clause writes them (cast).

\begin{stmt}
\sit{relation}{x\ \text{zlvl}\ L:\iff x,L\in \N,\ x \le 2^L,\ 2^L \le x+2^L}
\end{stmt}

\begin{lemma}[{a relation at every level}]\label{lem:zlvlbase}
\begin{stmt}
\sit{core}{K\in \N:\ 2^K\ \text{zlvl}\ K}
\end{stmt}
\end{lemma}

\begin{lemma}[{an equation of natural numbers substituted}]\label{lem:zsubst}
\begin{stmt}
\sit{right}{D\in \N:\ \forall L\in \N:\ (\exists L'\in \N:\ L'+D = L)\Rightarrow \exists x\in \N:\ x\ \text{zlvl}\ L}
\sit{premise}{D,M\in \N:\ \forall L\in \N:\ L = M+D\Rightarrow \exists x\in \N:\ x\ \text{zlvl}\ L}
\sit{with}{D\in \N:\ \forall L\in \N:\ (\exists L'\in \N:\ L = L'+D)\Rightarrow \exists x\in \N:\ x\ \text{zlvl}\ L}
\sit{rule}{D\in \N:\ (\forall L,x\in \N:\ x\ \text{zlvl}\ L\Rightarrow 2 x\ \text{zlvl}\ L+1)\Rightarrow \forall L\in \N:\ (\exists L'\in \N:\ L = L'+D+1)\Rightarrow \exists x\in \N:\ x\ \text{zlvl}\ L}
\sit{cast}{D\in \N,\ q\in \Q_{\ge0}:\ \forall L\in \N:\ (\exists L'\in \N:\ L = L'+D,\ q \le L')\Rightarrow q+D \le L}
\end{stmt}
\end{lemma}

\begin{proof}
\begin{plines}
\pline{right}{By logic: \textup{(1)}~\(2^{L'+D}\ \text{zlvl}\ L'+D\) by \cref{lem:zlvlbase}(core). Take \(x=2^{L'+D}\). The goal follows from the facts at hand.}
\pline{premise}{By logic: \textup{(1)}~\(2^{M+D}\ \text{zlvl}\ M+D\) by \cref{lem:zlvlbase}(core). Take \(x=2^{M+D}\). The goal follows from the facts at hand.}
\pline{with}{\textup{(1)}~\(\exists L'\in \N:\ L = L'+D\) by logic. \textup{(2)}~\(2^{L'+D}\ \text{zlvl}\ L'+D\) by \cref{lem:zlvlbase}(core). Take \(x=2^{L'+D}\). The goal follows from the facts at hand.}
\pline{rule}{By logic: \textup{(1)}~\(2^{L'+D}\ \text{zlvl}\ L'+D\) by \cref{lem:zlvlbase}(core). By logic: take \(x=2\cdot 2^{L'+D}\). The goal follows from the facts at hand.}
\pline{cast}{By logic: \textup{(1)}~\(q+D \le L'+D\) from the facts at hand. The goal follows by logic.}
\end{plines}
\end{proof}

\begin{lemma}[{a subset of regions by logic}]\label{lem:zsubset}
logic on a goal A subseteq B: the point w and its membership in A introduced, the goal its membership in B (a
union read as the disjunction of the memberships); the rules read the membership, and the word after them (with, a
case split on the point by cases) names the point w. The parts of a premise that logic destructs are rules as the
hypotheses are (conj).
\begin{stmt}
\sit{inc}{n\in \N,\ K\in \mathrm{Loss}(n),\ R,S\in \mathrm{Region}(n),\ (\forall y\in \mathrm{Point}(n):\ (\exists v\in R:\ \forall i:\ v_i = y_i)\iff K(y) \le 1),\ (\forall y\in \mathrm{Point}(n):\ (\exists v\in S:\ \forall i:\ v_i = y_i)\iff K(y) \le 2):\ R \subseteq S}
\sit{union}{n\in \N,\ K\in \mathrm{Loss}(n),\ R,S,U\in \mathrm{Region}(n),\ (\forall y\in \mathrm{Point}(n):\ (\exists v\in R:\ \forall i:\ v_i = y_i)\iff K(y) \le 3),\ (\forall y\in \mathrm{Point}(n):\ (\exists v\in S:\ \forall i:\ v_i = y_i)\iff K(y) \le 1),\ (\forall y\in \mathrm{Point}(n):\ (\exists v\in U:\ \forall i:\ v_i = y_i)\iff 1 \le K(y)):\ R \subseteq S\cup U}
\sit{conj}{n\in \N,\ K\in \mathrm{Loss}(n),\ R,S\in \mathrm{Region}(n):\ (\forall y\in R:\ K(y) \le 1),\ (\forall y\in \mathrm{Point}(n):\ K(y) \le 2\Rightarrow \exists v\in S:\ \forall i:\ v_i = y_i)\Rightarrow R \subseteq S}
\end{stmt}
\end{lemma}

\begin{proof}
\begin{plines}
\pline{inc}{By logic: \textup{(1)}~\(K(w) \le 2\) from the facts at hand. The goal follows by logic.}
\pline{union}{By logic: Case \((\mathrm{c}_{1})\) \(K(w) \le 1\): the goal follows by logic. Case \((\bar{\mathrm{c}}_{1})\) \(K(w) > 1\): \textup{(1)}~\(1 \le K(w)\) from the facts at hand. The goal follows by logic.}
\pline{conj}{By logic: \textup{(1)}~\(K(w) \le 2\) from the facts at hand. The goal follows by logic.}
\end{plines}
\end{proof}

\begin{lemma}[{a universal goal on the name of a witness}]\label{lem:zfresh}
A universal goal on a name in scope: after with destructs exists i : P, a goal forall i : Q introduces its index as
i_2 (as logic names a clashing witness), so that the facts on the witness i keep their meaning and their Lean text.
\begin{stmt}
\sit{wit}{n\in \N,\ a\in \mathrm{Grid}(n),\ (\exists i:\ 2 \le a_i):\ (\exists j:\ 1 \le a_j),\ \forall i:\ 1 \le a_i\Rightarrow 2 \le a_i+a_i}
\sit{rule}{n\in \N,\ a,b\in \mathrm{Grid}(n),\ (\exists i:\ 1 \le a_i),\ (\forall i:\ 1 \le a_i\Rightarrow 1 \le b_i):\ (\exists j:\ 1 \le a_j),\ \forall i:\ 1 \le a_i\Rightarrow 2 \le b_i+b_i}
\end{stmt}
\end{lemma}

\begin{proof}
\begin{plines}
\pline{wit}{\textup{(1)}~\(\exists i:\ 2 \le a_i\) by logic. Part (i): take \(j=i\). The goal follows from the facts at hand. Part (ii): the goal follows from the facts at hand.}
\pline{rule}{\textup{(1)}~\(\exists i:\ 1 \le a_i\) by logic. Part (i): take \(j=i\). The goal follows from the facts at hand. Part (ii): the goal follows by logic and from the facts at hand.}
\end{plines}
\end{proof}

\begin{lemma}[{the reading of a program row with a guard}]\label{lem:zprogguard}
\begin{stmt}
\sit{given}{\mathrm{kq},\eta,z\in \Q_{\ge0},\ 0 < \mathrm{kq},\ (1+\eta) z \ge 1}
\sit{(1)}{\mathrm{kq}/(1+\eta) \le \mathrm{kq}\cdot z}
\end{stmt}
\end{lemma}

\begin{proof}
\begin{plines}
\pline{(1)}{The goal follows by \cref{lem:progstrict}(1) (at \(x=z\)).}
\end{plines}
\end{proof}

\begin{lemma}[{a program as one object}]\label{lem:zbundle}
A bundle: a relation on it as a hypothesis, read by its fields; inside a bracket, decided classically; the
word unfold reads the relations of a goal by their fields (R_iff), so that two sides that read the same once
unfolded are equal.
\begin{stmt}
\sit{hyp}{k,n,L\in \N,\ \mathrm{Pr}\in \mathrm{Prog}(k,n),\ s\in \mathrm{Grid}(n),\ s\ \text{zfeasible}\ \mathrm{Pr}\ \text{zlevel}\ L:\ \forall \iota:\ \textstyle\sum_i \mathrm{Pr}.P_{\iota,i}\cdot s_i+\mathrm{Pr}.\mathrm{op}_{\iota} \ge \mathrm{Pr}.\mathrm{om}_{\iota}}
\sit{count}{k,n,L\in \N,\ \mathrm{Pr}\in \mathrm{Prog}(k,n),\ b\in \mathrm{Grid}(n):\ \textstyle\sum_{s\le b} \mathbf 1[s\ \text{zfeasible}\ \mathrm{Pr}\ \text{zlevel}\ L] \le \textstyle\sum_{s\le b} 1}
\sit{unfold}{k,n,L\in \N,\ \mathrm{Pr}\in \mathrm{Prog}(k,n),\ b\in \mathrm{Grid}(n):\ \textstyle\sum_{s\le b} \mathbf 1[s\ \text{zfeasible}\ \mathrm{Pr}\ \text{zlevel}\ L] = \textstyle\sum_{s\le b} \mathbf 1[\forall \iota:\ \textstyle\sum_i \mathrm{Pr}.P_{\iota,i}\cdot s_i+\mathrm{Pr}.\mathrm{op}_{\iota} \ge \textstyle\sum_i \mathrm{Pr}.N_{\iota,i}\cdot s_i+\mathrm{Pr}.\mathrm{om}_{\iota}+\mathrm{Pr}.\mathrm{eps}_{\iota}\cdot L]}
\sit{build}{k,n,a\in \N:\ \exists \mathrm{Pr}\in \mathrm{Prog}(k,n):\ \forall \iota:\ (\forall i:\ \mathrm{Pr}.P_{\iota,i}+\mathrm{Pr}.N_{\iota,i} = 2),\ \mathrm{Pr}.\mathrm{op}_{\iota} = a+1,\ \mathrm{Pr}.\mathrm{om}_{\iota}+\mathrm{Pr}.\mathrm{eps}_{\iota} = 1}
\sit{cert}{k,a\in \N:\ \exists C\in \mathrm{ZCert}(k):\ \mathrm{C}.d = a+1,\ \forall \iota:\ \mathrm{C}.y_{\iota} \ge \mathrm{C}.d}
\sit{diag}{n,a\in \N:\ \exists \mathrm{Pr}\in \mathrm{Prog}(n,n):\ \forall \iota:\ (\forall i:\ \mathrm{Pr}.P_{\iota,i} = 2 \mathbf 1[\iota = i]),\ \textstyle\sum_i \mathrm{Pr}.P_{\iota,i} = 2}
\end{stmt}
\end{lemma}

\begin{proof}
\begin{plines}
\pline{hyp}{The goal follows by logic and from the facts at hand.}
\pline{count}{The goal follows term by term, by cases on the conditionals of the facts, then by computation in each case.}
\pline{unfold}{The goal follows by the definitions, from the facts at hand.}
\pline{build}{Take \(\mathrm{Pr}.P_{\iota,i}=2, \mathrm{Pr}.N_{\iota,i}=0, \mathrm{Pr}.\mathrm{op}_{\iota}=a+1, \mathrm{Pr}.\mathrm{om}_{\iota}=0, \mathrm{Pr}.\mathrm{eps}_{\iota}=1\). Each part of the goal follows from the facts at hand.}
\pline{cert}{Take \(\mathrm{C}.y_{\iota}=a+2, \mathrm{C}.d=a+1\). Each part of the goal follows from the facts at hand.}
\pline{diag}{Take \(\mathrm{Pr}.P_{\iota,i}=2 \mathbf 1[\iota = i], \mathrm{Pr}.N_{\iota,i}=0, \mathrm{Pr}.\mathrm{op}_{\iota}=a+1, \mathrm{Pr}.\mathrm{om}_{\iota}=0, \mathrm{Pr}.\mathrm{eps}_{\iota}=1\). The parts of the goal follow: (i) from the facts at hand; (ii) term by term.}
\end{plines}
\end{proof}

\begin{lemma}[{two indices compare}]\label{lem:zindex}
Indices are counts, with their order and equality: two indices of one range compare, in an indicator or as a
guard; the indicator of equality picks the coordinate out of a sum.
\begin{stmt}
\sit{pick}{n\in \N,\ x\in \mathrm{Grid}(n):\ \forall i:\ \textstyle\sum_{\iota} \mathbf 1[\iota = i] x_{\iota} = x_i}
\sit{below}{n\in \N,\ x\in \mathrm{Grid}(n):\ \forall i:\ \textstyle\sum_{\iota} \mathbf 1[\iota < i] x_{\iota} \le \textstyle\sum_{\iota} \mathbf 1[\iota \le i] x_{\iota}}
\sit{guard}{n\in \N,\ x\in \mathrm{Grid}(n),\ (\forall i:\ \forall j:\ i \le j\Rightarrow x_i \le x_j):\ \forall i:\ \forall j:\ i < j\Rightarrow x_i \le x_j}
\end{stmt}
\end{lemma}

\begin{proof}
\begin{plines}
\pline{pick}{The goal follows term by term.}
\pline{below}{The goal follows term by term, by cases on the conditionals of the facts, then by computation in each case.}
\pline{guard}{The goal follows by logic and from the facts at hand.}
\end{plines}
\end{proof}

\begin{lemma}[{an index compares with an index of another range or with a count}]\label{lem:zcross}
An index compares with an index of another range, or with a count, as the count it is: the rows iota of a
tuple of tuples against the coordinates i of a row, an index up to a count t, the indicator of two indices of two
ranges in a sum, read term by term (no pick: the two ranges differ), and the rows of a second block, the row iota
at the coordinate i when iota = K + i.
\begin{stmt}
\sit{cross}{k,n\in \N,\ f\in \mathrm{Rows}(k,n),\ (\forall \iota:\ \forall i:\ \iota \le i\Rightarrow f_{\iota,i} \le 1):\ \forall \iota:\ \forall i:\ \iota < i\Rightarrow f_{\iota,i} \le 1}
\sit{count}{n,t\in \N,\ x\in \mathrm{Grid}(n),\ (\forall i:\ i \le t\Rightarrow x_i \le 1):\ \forall i:\ i < t\Rightarrow x_i \le 1}
\sit{diag}{k,n\in \N,\ x\in \mathrm{Grid}(n),\ y\in \mathrm{Grid}(k):\ \forall \iota:\ \textstyle\sum_i \mathbf 1[\iota = i] y_{\iota} x_i \le y_{\iota}\cdot (\textstyle\sum_i x_i)}
\sit{block}{k,n,K\in \N,\ f\in \mathrm{Rows}(k,n),\ (\forall \iota:\ \forall i:\ \iota = K+i\Rightarrow f_{\iota,i} \le 1):\ \forall \iota:\ \forall i:\ K+i = \iota\Rightarrow f_{\iota,i} \le 1}
\end{stmt}
\end{lemma}

\begin{proof}
\begin{plines}
\pline{cross}{The goal follows by logic and from the facts at hand.}
\pline{count}{The goal follows by logic and from the facts at hand.}
\pline{diag}{The goal follows term by term, by cases on the conditionals of the facts, then by computation in each case.}
\pline{block}{The goal follows by logic and from the facts at hand.}
\end{plines}
\end{proof}

A relation used at other names of its dimensions: the dimensions of the use are read from the declarations of its
arguments (k, n := n, k, at once), so that its field n <= t reads k <= t at the use, and casts to the goal over QQ.

\begin{stmt}
\sit{relation}{f\ \text{zwide}\ t:\iff k,n\in \N,\ f\in \mathrm{Rows}(k,n),\ t\in \N,\ n \le t}
\end{stmt}

\begin{lemma}[{a relation at other names of its dimensions}]\label{lem:zwide}
\begin{stmt}
\sit{swap}{k,n,t\in \N,\ x\in \Q_{\ge0},\ g\in \mathrm{Rows}(n,k),\ g\ \text{zwide}\ t,\ x \le k:\ x \le t}
\end{stmt}
\end{lemma}

\begin{proof}
\begin{plines}
\pline{swap}{The goal follows from the facts at hand.}
\end{plines}
\end{proof}

A bundle that take builds part by part, at other names of its dimensions: the fields of the relation read them from
the set of the existential (r, h for k, n), as for a declared argument.

\begin{stmt}
\sit{type}{\mathrm{ZBox}(k,n):=\{\mathrm{lo},\mathrm{hi}\in \mathrm{Grid}(n),\ w\in \mathrm{Grid}(k),\ d\in \N\}}
\sit{relation}{B\ \text{zfit}\ t:\iff k,n\in \N,\ B\in \mathrm{ZBox}(k,n),\ t\in \N,\ n = \mathrm{B}.d+1,\ k \le t}
\end{stmt}

\begin{lemma}[{a bundle built by take at other names of its dimensions}]\label{lem:ztake}
\begin{stmt}
\sit{build}{r,h,t\in \N,\ h = 3,\ r \le t:\ \exists B\in \mathrm{ZBox}(r,h):\ B\ \text{zfit}\ t}
\end{stmt}
\end{lemma}

\begin{proof}
\begin{plines}
\pline{build}{Take \(\mathrm{B}.\mathrm{lo}_i=0, \mathrm{B}.\mathrm{hi}_i=1, \mathrm{B}.w_j=0, \mathrm{B}.d=2\). Each part of the goal follows from the facts at hand.}
\end{plines}
\end{proof}

\begin{lemma}[{one reference for a conjunction}]\label{lem:zpair}
One reference proves a conjunction whole, as it proves any goal: its conclusion is the conjunction.
\begin{stmt}
\sit{core both}{x,y\in \Q_{\ge0},\ x \le y:\ x \le y+1,\ x \le y+2}
\sit{pair}{x,y\in \Q_{\ge0},\ x \le y:\ x \le y+1,\ x \le y+2}
\end{stmt}
\end{lemma}

\begin{proof}
\begin{plines}
\pline{pair}{The goal follows by \cref{lem:zpair}(core both).}
\end{plines}
\end{proof}

\begin{lemma}[{a loss given by a formula}]\label{lem:zpointfun}
A function on points given by a formula in its point (`K(v) := t`, as a family by a formula in its index): the rule
of a loss at the loss that sums the coordinates, its hypothesis and its conclusion read at the formula.
\begin{stmt}
\sit{base}{n\in \N,\ K\in \mathrm{Loss}(n),\ w\in \mathrm{Point}(n),\ a\in \Q_{\ge0},\ K(w) \le a:\ K(w) \le a+1}
\sit{inst}{n\in \N,\ w\in \mathrm{Point}(n),\ a\in \Q_{\ge0},\ \textstyle\sum_i w_i \le a:\ \textstyle\sum_i w_i \le a+1}
\end{stmt}
\end{lemma}

\begin{proof}
\begin{plines}
\pline{base}{The goal follows from the facts at hand.}
\pline{inst}{The goal follows by \cref{lem:zpointfun}(base).}
\end{plines}
\end{proof}

\begin{stmt}
\sit{type}{\mathrm{Mesh}(n):=\mathrm{Grid}(n)\to \Q_{\ge0}}
\end{stmt}

\begin{lemma}[{a family of functions given by a formula}]\label{lem:zfamfun}
A family of functions on points given by a formula in its index and its point (`W_(p)(s) := t`, as `K(v) := t` for
one function): the sum of the bounds of a family of weights on a box, at weights 2^(-s) scaled by a family of masses.
\begin{stmt}
\sit{base}{P,n\in \N,\ b\in \mathrm{Grid}(n),\ W_1,\dots,W_P\in \mathrm{Mesh}(n),\ U_1,\dots,U_P\in \Q_{\ge0},\ (\forall p\in \N_{\ge1}:\ p \le P\Rightarrow \textstyle\sum_{s\le b} W_p(s) \le U_p):\ \textstyle\sum_{p=1}^P (\textstyle\sum_{s\le b} W_p(s)) \le \textstyle\sum_{p=1}^P U_p}
\sit{inst}{P,n\in \N,\ b\in \mathrm{Grid}(n),\ a_1,\dots,a_P\in \Q_{\ge0},\ (\forall p\in \N_{\ge1}:\ p \le P\Rightarrow \textstyle\sum_{s\le b} a_p/2^{\textstyle\sum_i s_i} \le a_p+1):\ \textstyle\sum_{p=1}^P (\textstyle\sum_{s\le b} a_p/2^{\textstyle\sum_i s_i}) \le \textstyle\sum_{p=1}^P (a_p+1)}
\end{stmt}
\end{lemma}

\begin{proof}
\begin{plines}
\pline{base}{The goal follows term by term, by logic.}
\pline{inst}{The goal follows by \cref{lem:zfamfun}(base).}
\end{plines}
\end{proof}

\begin{lemma}[{an index read as a count}]\label{lem:zidxcount}
An index over the coordinates read as the count it is in a term, as the coefficient 2 * iota of a
row: a count read from an index indexes nothing.
\end{lemma}

\begin{lemma}[{the bound of an index and the guard of a name, by name}]\label{lem:zidxbound}
The bound of an index and the guard of a name, each by name. An index over the coordinates of the tuples of length
n is a count from 1 to n: its bound i <= n is a fact of the step (Lean: i.isLt), which a certificate names hb(i)
and a case reads by name; the guard of a name of the given, k <= x for x in NN>=k, 0 < x for x in QQ>0, is named
hp(x). An index above the last index but one is the last: a value from an inner and an outer bound.
\end{lemma}

\begin{lemma}[{the tie of a denominator, by name}]\label{lem:zdentie}
The tie of a denominator. A certificate whose multipliers or facts are quotients is exact once the ties d * d^-1 = 1
of its denominators are added, each denominator above 0 by the signs of its names; a mass substituted into a program
gives the fact of its denominator too, so that the identity of the value reads its tie.
\begin{stmt}
\sit{cell}{A,B,D\in \N_{\ge1},\ A = B+D:\ B/A+D/A = 1}
\sit{mass}{a,b,d\in \Q_{>0},\ a = b+d:\ b/a+d/a = 1}
\sit{floor}{A,T\in \N_{\ge1},\ A \le T:\ 1 \le T/A}
\sit{sub}{t,\alpha\in \Q_{>0},\ \alpha = 2 t/(2+2 t):\ \operatorname{prog}(v\in \Q_{\ge0},\ v \ge 1:\ \alpha v) = (2 t/(2+2 t),1)}
\end{stmt}
\end{lemma}

\begin{proof}
\begin{plines}
\pline{cell}{The goal follows from \(1/A\) times \((\mathrm{H}_{1})\).}
\pline{mass}{The goal follows from \(1/a\) times \((\mathrm{H}_{1})\).}
\pline{floor}{The goal follows from \(1/A\) times \((\mathrm{H}_{1})\).}
\pline{sub}{The goal follows by the prices \(t/(1+t)\) of the rows at the points \(v=1\).}
\end{plines}
\end{proof}

\begin{lemma}[{the signs and the ties of a certificate, by name}]\label{lem:zcertsign}
The signs and the ties of a certificate, by name. A sum is above 0 when one of its parts is and the others
are at least 0 (a count plus one, as a denominator or as a multiplier of a strict fact). The inverse of a sum is
read with its tie d * d^-1 = 1, for a sum of masses as for a count plus one. A multiplier that cancels a quotient
of a fact reads the tie of its denominator, and so does an identity of quotients whose two sides are one element of
the Laurent ring. A product of two strict facts is strict.
\end{lemma}

\begin{lemma}[{a reference read by rule}]\label{lem:zrefread}
A reference read by rule: its conclusion at the data is the goal, a conjunct of it, or the goal with its two
sides reversed, and its proof is the term of the reference, with its projection or `.symm`.
\end{lemma}

\begin{lemma}[{a least of a least, at its branch}]\label{lem:znestmin}
A least whose term is a least, nested or written flat (min(a, b, c) is min(a, min(b, c)), as Lean writes it), read
at its branch by a fact of the step that names both of its terms: the tie min(a, min(b, c)) = a by
a <= min(b, c), no case split.
\begin{stmt}
\sit{nested}{a,b,c,y\in \Q_{\ge0},\ y \ge \min(a,\min(b,c)),\ b \le c,\ a \le b:\ y \ge a}
\sit{flat}{a,b,c,y\in \Q_{\ge0},\ y \ge \min(a,b,c),\ b \le c,\ a \le b:\ y \ge a}
\end{stmt}
\end{lemma}

\begin{proof}
\begin{plines}
\pline{nested}{\textup{(1)}~\(a \le \min(b,c)\) from \((\mathrm{H}_{7})\). The goal follows from \((\mathrm{H}_{5})\).}
\pline{flat}{\textup{(1)}~\(a \le \min(b,c)\) from \((\mathrm{H}_{7})\). The goal follows from \((\mathrm{H}_{5})\).}
\end{plines}
\end{proof}

\begin{stmt}
\sit{type}{\mathrm{Starz}(k):=\{N_0,w_0\in \N_{\ge1},\ N,w\in (\N_{\ge1})^{k}\}}
\sit{relation}{\mathrm{St}\ \text{zplane}\ s:\iff k,s\in \N_{\ge1},\ \mathrm{St}\in \mathrm{Starz}(k),\ \mathrm{St}.N_0\cdot s = \textstyle\sum_j \mathrm{St}.N_j,\ \mathrm{St}.w_0\cdot s+k = 2+(\textstyle\sum_j \mathrm{St}.w_j)}
\end{stmt}

\begin{lemma}[{the ties of a certificate at its branches, and the premises of a step}]\label{lem:zfollow}
The ties of a certificate read at its branches: the sides of a goal rewritten at the branch of its least, where the
branch brings a quotient (rewrite); a quotient by a count or a coordinate plus a positive constant, with the tie of
its inverse, and a rest at least 0 read with the denominator cleared (sumnat, sumpair); a least of a least decided
by a strict fact and a fact (strictmin); the leasts of an equation decided by the facts of a with (value); the
premise of an implication in the step of an induction, as its clause states it (walk); and the bound of the size of
a bundle as a field of a relation (plane).
\begin{stmt}
\sit{rewrite}{F\in \Q_{>0}:\ F \min(1/2,1/F) \le 1}
\sit{sumnat}{b\in \Q_{>0},\ \eta\in \N,\ b \ge 2:\ b/2 > 1/2+\max(b-1,0)/(3+\eta)}
\sit{sumpair}{b\in \Q_{>0},\ \eta\in \Q_{\ge0},\ b \ge 1:\ b/2 \ge 1/2+\max(b-1,0)/(3+\eta)}
\sit{strictmin}{a,b,c,s\in \Q_{\ge0},\ s \ge \min(c,\min(a,b)),\ a \le b,\ \min(a,b) < c:\ s \ge a}
\sit{value}{k',\eta\in \Q_{\ge0}:\ \min(1+(1+k')/(3+\eta),(3+k')/3) = 1+\min(k'/3,(1+k')/(3+\eta))}
\sit{(walk)}{N\in \N,\ n\in \N_{\ge1},\ p_1,\dots,p_{N+1}\in \Q_{\ge0},\ (\forall j\in \N_{\ge1}:\ j \le N\Rightarrow p_j \le p_{j+1}):\ n \le N+1\Rightarrow p_{1} \le p_n}
\sit{plane}{k,s\in \N_{\ge1},\ \mathrm{St}\in \mathrm{Starz}(k),\ \mathrm{St}\ \text{zplane}\ s:\ 1 \le k}
\end{stmt}
\end{lemma}

\begin{proof}
\begin{plines}
\pline{rewrite}{Case \((\mathrm{c}_{1})\) \(F \le 2\): \textup{(1)}~\(1/2 \le 1/F\) from \(1/(2 F)\) times \((\mathrm{c}_{1})\). The goal follows from \(1/2\) times \((\mathrm{c}_{1})\). Case \((\bar{\mathrm{c}}_{1})\) \(F > 2\): \textup{(2)}~\(1/F \le 1/2\) from \(1/(2 F)\) times \((\bar{\mathrm{c}}_{1})\). The goal follows by computation.}
\pline{sumnat}{\textup{(1)}~\(\max(b-1,0) = b-1\) by computation. The goal follows combining \((1+\eta)/(2\cdot (3+\eta))\) times \((\mathrm{H}_{1})\), with \(1/(3+\eta)\) times \textup{(1)} read from right to left.}
\pline{sumpair}{\textup{(1)}~\(\max(b-1,0) = b-1\) by computation. The goal follows combining \((1+\eta)/(2\cdot (3+\eta))\) times \((\mathrm{H}_{2})\), with \(1/(3+\eta)\) times \textup{(1)} read from right to left.}
\pline{strictmin}{The goal follows from \((\mathrm{H}_{5})\).}
\pline{value}{Case \((\mathrm{c}_{1})\) \(k'/3 \le (1+k')/(3+\eta)\): \textup{(1)}~\((3+k')/3 \le 1+(1+k')/(3+\eta)\) from \((\mathrm{c}_{1})\). The goal follows by computation. Case \((\bar{\mathrm{c}}_{1})\) \(k'/3 > (1+k')/(3+\eta)\): \textup{(2)}~\((1+k')/(3+\eta) \le k'/3\) from \((\bar{\mathrm{c}}_{1})\). \textup{(3)}~\(1+(1+k')/(3+\eta) \le (3+k')/3\) from \textup{(2)}. The goal follows by computation.}
\pline{(walk)}{By induction on \(n\). Base: the goal follows from the facts at hand. Step: \textup{(1)}~\(n \le N\) from \((\mathrm{I}_{2})\). \textup{(2)}~\(n \le N+1\) from the facts at hand. \textup{(3)}~\(p_{1} \le p_n\) by logic. \textup{(4)}~\(p_n \le p_{n+1}\) by logic. The goal follows combining \textup{(3)} and \textup{(4)}.}
\pline{plane}{The goal follows by logic.}
\end{plines}
\end{proof}

\begin{lemma}[{the inverse of a count, and the sides of a product of facts}]\label{lem:zinvcount}
The inverse of a count, and a product of facts with its sides: a strict goal through the tie
pi_(i) = 1 / E_(i), with the sign of 1 / E_(i) by the fact 1 <= E_(X) at its index (face, span); and a product of
two facts times a quotient by a count plus one, whose ties are read from the sides of the product (product).
\begin{stmt}
\sit{face}{n\in \N,\ E\in \mathrm{Grid}(n),\ \pi\in \mathrm{Point}(n),\ (\forall j:\ \pi_j = 1/E_j),\ (\forall X:\ 1 \le E_X):\ \forall i:\ 0 < \pi_i}
\sit{span}{n\in \N,\ E\in \mathrm{Grid}(n),\ \pi\in \mathrm{Point}(n),\ (\forall j:\ \pi_j = 1/E_j),\ (\forall X:\ 1 \le E_X):\ \forall i:\ 0 < \pi_i+\pi_i}
\sit{product}{J\in \N,\ k\in \Q_{>0},\ c,x,w\in \Q_{\ge0},\ v_1,\dots,v_J\in \Q_{\ge0},\ J\cdot (J+1) \le 2 k,\ 2 k \le (J+1)\cdot (J+2),\ 2 c \ge 1,\ 2\cdot (J+1) x+2 w \ge 1,\ \textstyle\sum_{i=1}^J v_i+J\cdot (J+1)/2 x \ge J/2:\ c+k x+(\textstyle\sum_{i=1}^J v_i)+w \ge (J+2)/4+k/(2\cdot (J+1))}
\end{stmt}
\end{lemma}

\begin{proof}
\begin{plines}
\pline{face}{By logic: \textup{(1)}~\(\pi_i = 1/E_i\) by logic. \textup{(2)}~\(0 < 1/E_i\) from the facts at hand. The goal follows combining \textup{(2)}, with \textup{(1)} read from right to left.}
\pline{span}{By logic: \textup{(1)}~\(\pi_i = 1/E_i\) by logic. \textup{(2)}~\(0 < 1/E_i\) from the facts at hand. The goal follows combining \(2\) times \textup{(2)}, with \(2\) times \textup{(1)} read from right to left.}
\pline{product}{The goal follows combining \(1/2\) times \((\mathrm{H}_{7})\), \((\mathrm{H}_{9})\), \(1/(4\cdot (J+1))\) times the product of \((\mathrm{H}_{5})\) and \((\mathrm{H}_{8})\) and \(w/(2\cdot (J+1))\) times \((\mathrm{H}_{6})\).}
\end{plines}
\end{proof}

\begin{lemma}[{a power of a count as a coefficient, and a fact at 0}]\label{lem:zpowcoef}
A power of a count as a coefficient of a certificate, read with its cast pushed in as the goal reads it; and a
universal fact at an index 0, its index term reduced before the cast.
\begin{stmt}
\sit{inst}{N,\nu,p\in \N,\ c\in \mathrm{Line},\ (\forall L\in \N:\ c(L+N) p^{\nu} = c(L)):\ c(N) p^{\nu} = c(0)}
\end{stmt}
\end{lemma}

\begin{proof}
\begin{plines}
\pline{inst}{The goal follows by \((\mathrm{H}_{1})\) at \(L=0\).}
\end{plines}
\end{proof}

\begin{lemma}[{a fact at a given point}]\label{lem:zfactat}
A fact of the step at a given term: a universal over the points of Grid(n) read at the point whose coordinate l
is iverson(l = j), each sum of a term times the indicator of its index collapsed to the term at j.
\end{lemma}

\begin{lemma}[{the laws of sums, one at each place}]\label{lem:zsumlaws}
An equation of terms with sums by `eq`, each place where the sides differ given by one law: an identity of rings
over the same sums, a sum of a sum split, a sum against the indicator of its index collapsed, a factor moved into
a sum, two sums exchanged, a factor moved out of a sum; the places reached through the sums and the terms.
\end{lemma}

\begin{lemma}[{a fact read at the names of a place}]\label{lem:zfactplace}
The place where two sides differ given by a fact over two indices, and by a fact over an index and a point, each
read at the names of the place by matching its two sides.
\end{lemma}

\begin{lemma}[{the halves of the bideterminant, the two sides of a tie}]\label{lem:zbidet}
The bideterminant: bidetP(T, U) and bidetM(T, U), the two halves of the determinant of a couple of square
tables (T, U) standing for T - U, over the semiring of couples (Akian, Gaubert and Guterman, arXiv:0812.3496,
Def. 4.14); of one table, bidetP(A) and bidetM(A), the bideterminant (|A|^+, |A|^-) of Gondran and Minoux, the sums
over the even and over the odd permutations, so that det A = |A|^+ - |A|^- over a ring. A half is an atom of a
certificate, never unfolded, and a tie, a comparison or a deficit reads each couple only through T - U: the
difference of its sides is invariant under one mass added to both halves of a couple, as in the tie of crossed
products of two couples, the form of the law of Binet and Cauchy (core cross).
\end{lemma}

\begin{lemma}[{the cast of a product of counts}]\label{lem:zcountcast}
The certificate of an order step reads the cast of a sum or a product of counts as the sum or the product of
their casts.
\end{lemma}

\begin{lemma}[{the integer rounding of a strict order of counts}]\label{lem:zsucc}
A strict order of counts a < b is a + 1 <= b (the strict order of counts is the order of the successor): where
the strict reading leaves the data of V no certificate, the strict orders of counts that they name are read so.
\end{lemma}

\begin{lemma}[{a multiplier of several terms, its sign by name}]\label{lem:zmultsign}
An order step whose multiplier has several terms: the inequality scaled by it, its sign written term by term
by name (a count by its cast, a sum of counts term by term); an equation times such a multiplier as it is.
\end{lemma}

\begin{lemma}[{the guards of binders by name}]\label{lem:zguardname}
The guards of the binders by name: a goal that is the bound of a member of a family of counts or the sign of a
member of a family of masses is closed by the guard of the family at the member; a reference that reads a family
whole takes the family's guard for its own (hp_e), the strict sign weakened where it asks for 0 <= e_(i).
\end{lemma}

\begin{lemma}[{a deficit of counts as a witness}]\label{lem:zdeficitwit}
A witness that is the deficit max(p - 1, 0) of two counts: the certificate reads it as p - 1 by the bound
1 <= p that a fact states (max_eq_left, Nat.sub_add_cancel), with no search on the truncated difference; a case
whose negation is pushed in by one proof term.
\end{lemma}

\begin{lemma}[{the rest of a certificate signed by name}]\label{lem:zrestsign}
\end{lemma}

\begin{lemma}[{the conditionals of a goal decided by the facts of the step}]\label{lem:zcondname}
\end{lemma}

\begin{lemma}[{the laws of a range by name}]\label{lem:zrangelaw}
\end{lemma}

\begin{lemma}[{the shift of an index, a term free of the index, a smaller range}]\label{lem:zrangeshift}
The index of a range shifted by an offset free of it, either side, a term free of the index summed over a range (the
count of the range times it), and a sum over a range at most the sum of the same terms over a larger range, the
terms at least 0.
\end{lemma}

\begin{lemma}[{the sum over a box of a product over the coordinates}]\label{lem:zboxprod}
The law of a box: a sum over the lattice points of boxes of a product over the coordinates is the product
over the coordinates of the sums over their ranges, for one box, for two boxes, either side, and at each value of
a shared scale with a factor free of the points.
\end{lemma}

\begin{lemma}[{the strict sign of a sum over a nonempty range by name}]\label{lem:zsumsign}
\end{lemma}

\begin{lemma}[{a bound of counts by name}]\label{lem:znatbound}
\end{lemma}

\begin{lemma}[{a universal goal by a reference whole}]\label{lem:zwholeref}
\end{lemma}

\begin{lemma}[{a deficit of two sums read by the bound that names them}]\label{lem:zdeficitsum}
\end{lemma}

\begin{lemma}[{a rule read at the point that the goal names}]\label{lem:zgoalpoint}
\end{lemma}

\begin{lemma}[{a family on the counts by a formula}]\label{lem:zfamformula}
\end{lemma}

\begin{lemma}[{a family by its formula meets its defining hypothesis}]\label{lem:zfamdef}
\end{lemma}

\begin{lemma}[{the goal as an instance of a fact, by name}]\label{lem:zrulenamed}
\end{lemma}

\begin{lemma}[{a disjunction resolved by the negation of a case}]\label{lem:zdisjnamed}
\end{lemma}

\begin{lemma}[{a fact named up to the name its sum binds, a bound read either way round, two facts that clash}]\label{lem:zfactkey}
\end{lemma}

\begin{lemma}[{tuples of a size n + n'}]\label{lem:wsize}
Tuples of a size n + n' of two counts, read by the Mathlib reading at the sizes n and n' on its n + n'
coordinates: a loss at a joined point, a point, a sum and a box of that size.
\end{lemma}

\begin{lemma}[{a tuple binder on a positive set}]\label{lem:ztuplebind}
A tuple written in its binder, x in X^n or x in (X^n)^k with X one of QQ>0, QQ>=0, NN: the binder
reads as the type row it unfolds to, over Mathlib and over the kernel (Weight(n) := (QQ>0)^n is the same binder).
\end{lemma}

\begin{lemma}[{a row over the positive masses and counts}]\label{lem:zpositive}
A row whose names are all in QQ>0 or NN>=1 (a tuple of them), its terms built by +, *, / and powers from them and
from numerals at least 1: over the kernel it is printed over Mass and Count themselves, with no
Mass0.of and no Count0 in its terms; the size of a tuple stays a Count0 (Idx is on Count0). A row with one name
that holds zero is printed over the zero layer (core zero).
\end{lemma}

The positive counts as coordinates, and functions with positive values: a tuple of NN>=1, a
tuple of such tuples, a function to QQ>0 or to NN>=1, and a part of a bundle in NN>=1 (a tuple of them, or one),
each bound or taken as a witness. Over the kernel a coordinate in NN>=1 is a Count and a value in QQ>0 is a Mass;
the bound of each (at least 1, positive) is read by the bridge and is not a hypothesis of the row.

\begin{stmt}
\sit{type}{\mathrm{PosGrid}(n):=(\N_{\ge1})^{n}}
\sit{type}{\mathrm{PosRows}(k,n):=((\N_{\ge1})^{n})^{k}}
\sit{type}{\mathrm{PosLine}:=\N\to \Q_{>0}}
\sit{type}{\mathrm{StepLine}:=\N\to \N_{\ge1}}
\sit{type}{\mathrm{PosValues}(n):=\mathrm{Grid}(n)\to \Q_{>0}}
\sit{type}{\mathrm{PosCert}(k):=\{e\in \mathrm{PosGrid}(k),\ m\in \N_{\ge1},\ d\in \N\}}
\end{stmt}

\begin{lemma}[{tuples of positive counts and functions with positive values}]\label{lem:zposcount}
\begin{stmt}
\sit{core take part}{k\in \N,\ B\in \mathrm{PosCert}(k):\ \exists A\in \mathrm{Weight}(k):\ \forall i:\ A_i = \mathrm{B}.m\cdot \mathrm{B}.e_i}
\end{stmt}
\end{lemma}

\begin{proof}
\begin{plines}
\pline{core take part}{Take \(A_i=\mathrm{B}.m\cdot \mathrm{B}.e_i\). The goal follows from the facts at hand.}
\end{plines}
\end{proof}
